%%%%%%%%%%%%%%%%%%%%%%%%%%%%%%%%%%%%%%%%%%%%%%%%%%%%%%%%%%%%%%%%%%%%%%%%%%%%%%
%% UniReMoE manuscript -- OUP authoring template (Bioinformatics-style)
%% Build: ./build.sh (Tectonic or TeX Live); see README.md.
%%
%% PLACEHOLDER POLICY
%%   \NA{}   -- a result cell whose numerical evidence has not yet been inserted
%%   \TODO{} -- prose, citation or metadata the authors must supply
%%   \figph  -- a framed box standing in for a figure that is not drawn yet
%% All three helpers must be deleted before submission.
%%%%%%%%%%%%%%%%%%%%%%%%%%%%%%%%%%%%%%%%%%%%%%%%%%%%%%%%%%%%%%%%%%%%%%%%%%%%%%
\documentclass[webpdf,modern,large,namedate]{oup-authoring-template}

\graphicspath{{Fig/}}

% OUP's print page-box hook emits PostScript pdfmark commands. XeTeX/Tectonic
% uses xdvipdfmx; use its paper-size special for this web-PDF manuscript.
\ifdefined\XeTeXversion
  \makeatletter
  \def\shipout@PageObjects{%
    \special{papersize=\the\paperwidth,\the\paperheight}}
  \let\rest@dvi@pages\shipout@PageObjects
  \makeatother
\fi

% The class uses \societylogo{} in the page header but never defines it.
\newcommand{\societylogo}{}

% Theorem-like environments: the class supplies the styles, not the environments.
\theoremstyle{thmstyleone}
\newtheorem{theorem}{Theorem}
\newtheorem{proposition}[theorem]{Proposition}
\theoremstyle{thmstyletwo}
\newtheorem{example}{Example}
\newtheorem{remark}{Remark}
\theoremstyle{thmstylethree}
\newtheorem{definition}{Definition}

%%%%%%%%%%%%%%%%%%%%%%%%%%%%%%%%%%%%%%%%%%%%%%%%%%%%%%%%%%%%%%%%%%%%%%%%%%%%%%
%% Placeholder helpers -- DELETE BEFORE SUBMISSION
%%%%%%%%%%%%%%%%%%%%%%%%%%%%%%%%%%%%%%%%%%%%%%%%%%%%%%%%%%%%%%%%%%%%%%%%%%%%%%
\DeclareRobustCommand{\figph}[3][\linewidth]{%
  \fbox{\parbox[c][#2][c]{\dimexpr #1-2\fboxsep-2\fboxrule\relax}{%
    \centering\small\textit{#3}}}}
\newcommand{\NA}{N/A}
%% Robust, because \TODO appears inside captions that are written to the .lof.
\DeclareRobustCommand{\TODO}[1]{\textcolor{red}{[TODO: #1]}}
%%%%%%%%%%%%%%%%%%%%%%%%%%%%%%%%%%%%%%%%%%%%%%%%%%%%%%%%%%%%%%%%%%%%%%%%%%%%%%

%% Stable notation. Terminology must not drift between sections.
%% \PH/\PT/\PHT are the PROSE forms and must stay plain expandable text: they
%% appear in \title[...] (which feeds \markboth) and in \subsection{...}, where
%% a \text-based definition would upset hyperref's PDF-string handling.
%% Use the upright \PHm/\PTm/\PHTm forms inside math mode instead -- the prose
%% forms would set as italic multi-letter variables, and \PHT's en dash would
%% become two binary minus signs in a subscript.
\newcommand{\method}{UniReMoE}
\newcommand{\PH}{pHLA}
\newcommand{\PT}{pTCR}
\newcommand{\PHT}{pHLA--TCR}
\newcommand{\PHm}{\mathrm{PH}}
\newcommand{\PTm}{\mathrm{PT}}
\newcommand{\PHTm}{\mathrm{PHT}}
\newcommand{\sg}{\mathrm{sg}}
\newcommand{\LN}{\mathrm{LN}}
\newcommand{\MHA}{\mathrm{MHA}}
\newcommand{\relu}{\mathrm{relu}}
\newcommand{\softmax}{\mathrm{softmax}}
\newcommand{\dmodel}{d}

\begin{document}

\journaltitle{Bioinformatics}
\DOI{DOI added during production}
\copyrightyear{2026}
\pubyear{2026}
\vol{XX}
\volume{XX}   % \vol feeds the running heads, \volume the first-page opener box
\issue{X}
\access{Published: Date added during production}
\appnotes{Original Paper}

\firstpage{1}

\subtitle{Sequence analysis}

%% Keep the optional short title. The class's no-argument \title{...} branch
%% calls \authormark, which the class never defines, so it errors out.
\title[Shared--private interaction experts for \PHT{} prediction]{%
\method{}: shared--private relation experts and label-aware consistency for
unified peptide--HLA--TCR binding prediction}

\author[1,$\ast$]{Junpeng Ma\ORCID{0000-0000-0000-0000}}
\author[1]{\TODO{co-author name}}
\author[2]{\TODO{co-author name}}

\address[1]{\orgdiv{\TODO{division}}, \orgname{\TODO{institution}},
\orgaddress{\street{\TODO{street}}, \postcode{\TODO{postcode}},
\state{\TODO{state}}, \country{\TODO{country}}}}
\address[2]{\orgdiv{\TODO{division}}, \orgname{\TODO{institution}},
\orgaddress{\street{\TODO{street}}, \postcode{\TODO{postcode}},
\state{\TODO{state}}, \country{\TODO{country}}}}

\corresp[$\ast$]{Corresponding author.
\href{mailto:corresponding@example.org}{corresponding@example.org}}

\received{Date}{0}{Year}
\revised{Date}{0}{Year}
\accepted{Date}{0}{Year}

\abstract{
\textbf{Motivation:} Peptide--HLA (\PH{}), peptide--TCR (\PT{}) and
peptide--HLA--TCR (\PHT{}) prediction ask related questions with different
inputs. Positive pair labels do not establish recognition in every HLA
context, although the absence of class \texttt{110} (both pairs positive,
triple negative) makes conjunction sufficient to reproduce the observed labels.
An explicit triple head avoids building this dataset-specific regularity into
the prediction rule: it can use the joint molecular context instead of being
restricted to a fixed combination of pair scores. Learning exceptions to that
regularity still requires discriminating supervision. The unified model must
also share useful features without forcing identical task transformations.\\
\textbf{Results:} We present \method{}, a sequence-based model with three
jointly supervised prediction heads over frozen protein-language-model residue
embeddings. The pair heads use only their respective inputs, while the triple
head models recognition of the peptide under the specified HLA. To accommodate
their different feature requirements, each interaction block uses a dense
mixture of two globally shared and two branch-private experts. Relation-specific
routers select a mixture per sample, and the experts operate at each peptide
residue. This adapts an established multi-task topology to interactions with
different input compositions and preserves their outputs for the triple branch.
An HLA-conditioned per-residue gate transfers independently supervised \PT{}
features into the triple branch and lets triple supervision flow back to the
\PT{} feature extractor. Ablation experiments show that adding this gated
connection improves both \PT{} and \PHT{} prediction. Because feature sharing
does not directly regulate the three output scores, we add four label-aware
directional penalties. They raise lagging pair scores on positive triples or
lower excessive triple scores when the corresponding pair and triple labels
are negative, without clipping all triple predictions to a hard pair-score bound.
\TODO{Insert the transfer-ablation metrics, effect sizes and evaluation
protocol; complete the other quantitative comparisons from their run records.}\\
\textbf{Availability and implementation:} Source code, configurations, the
data-split audit and trained models are available at \TODO{repository URL}.\\
\textbf{Contact:}
\href{mailto:corresponding@example.org}{corresponding@example.org}\\
\textbf{Supplementary information:} Supplementary data are available at
\textit{Bioinformatics} online.}

\keywords{peptide--HLA binding, TCR specificity, protein language model,
mixture of experts, multi-task learning, immunogenicity}
\keywords[Abbreviations]{PLM, protein language model; MoE, mixture of experts;
\PH{}, peptide--HLA; \PT{}, peptide--TCR; \PHT{}, peptide--HLA--TCR;
AUROC, area under the receiver operating characteristic curve;
AUPRC, area under the precision--recall curve;
MCC, Matthews correlation coefficient}

\maketitle

% Keep the expanded key points out of the fixed-height opening page.
\clearpage
\section*{Key Points}
\begin{itemize}
\item One sequence-based model serves \PH{}, \PT{} and \PHT{} with explicit
input boundaries. A separately supervised triple head avoids fixing its
decision to pairwise conjunction; verified \texttt{110} examples are needed
to test whether it learns the additional HLA-dependent distinction.
\item Shared and branch-private experts let each relation select how much
interaction capacity to reuse. Routing is sample-level, expert computation is
residue-level, and the transformed pair representations remain available to
the triple branch.
\item HLA-conditioned gated transfer supplies independently supervised \PT{}
features to \PHT{} and returns triple-task gradients to the \PT{} feature
extractor. Adding this connection improves both tasks in the ablation study.
\item Label-aware consistency addresses output-score disagreement left open
by representation sharing and transfer. It selects a label-supported correction
direction and applies a soft penalty, retaining separate predictions rather
than enforcing a hard probability hierarchy.
\end{itemize}



%%%%%%%%%%%%%%%%%%%%%%%%%%%%%%%%%%%%%%%%%%%%%%%%%%%%%%%%%%%%%%%%%%%%%%%%%%%%%%
\section{Introduction}\label{sec:intro}
%%%%%%%%%%%%%%%%%%%%%%%%%%%%%%%%%%%%%%%%%%%%%%%%%%%%%%%%%%%%%%%%%%%%%%%%%%%%%%

T cells recognise peptides displayed by human leukocyte antigen (HLA)
molecules on the cell surface. Two questions therefore precede a prediction
for a particular peptide--HLA--TCR combination: can this HLA bind the peptide,
and is the T-cell receptor (TCR) associated with recognition of that peptide?
The final question is more specific: can this TCR recognise the peptide as
presented by this HLA? HLA holds the peptide in a binding groove, while the TCR
contacts the exposed peptide together with parts of the HLA molecule
\citep{rudolph2006structural}. Changing the HLA can therefore change the
surface available to the receptor. These interactions are relevant to
neoantigen selection and TCR-based therapy, although binding alone does not
guarantee a cellular immune response \citep{schumacher2015neoantigens,
sahin2018personalized}.

We denote the three prediction tasks by peptide--HLA (\PH{}), peptide--TCR
(\PT{}) and peptide--HLA--TCR (\PHT{}). Established predictors address
\PH{} binding or presentation \citep{reynisson2020netmhcpan,
odonnell2020mhcflurry}, or score peptide--TCR compatibility without an explicit
HLA input \citep{springer2020ergo, weber2021titan}. The latter
input convention does not mean that a TCR recognises a free peptide in
isolation; it means that the presenting HLA is unspecified in that prediction.
A positive \PH{} pair and a positive \PT{} pair therefore need not establish
recognition of the specified triple. For example, a receptor associated with a
peptide in one HLA context need not recognise it under another HLA that can
also bind that peptide. Combining two predicted pair scores by a minimum or
product supplies a useful baseline, but does not explicitly model this
allele-conditioned recognition. A minimum is a heuristic; a product requires
additional factorisation assumptions to be interpreted as a joint probability.

The label coverage makes this distinction especially important. Write the
three binary labels in the order \PH{}, \PT{}, \PHT{}. Our corpus contains
\texttt{000}, \texttt{100}, \texttt{010} and \texttt{111}, but no
\texttt{110}: both pairs positive while the specified triple is negative.
Consequently, on every available fully labelled row,
$y_{\PHTm}=y_{\PHm}\land y_{\PTm}$. A rule that predicts a positive triple
whenever both pair labels are positive thus reproduces the observed triple
labels exactly. The \texttt{111} examples establish co-occurrence, not the
universal implication that two positive pairs are sufficient. This restricted
label distribution leaves a conjunction-based solution unchallenged. If a
system contains only pair heads and defines its triple decision by conjunction,
it turns that empirical regularity into a fixed prediction rule, with no
separate decision based on the specified HLA--peptide--TCR context.

This motivates retaining \PHT{} as a separately supervised task. Its head
operates on joint residue-level representations of all three entities, rather
than treating the two pair scores as sufficient to determine the answer.
The purpose is to avoid imposing pairwise sufficiency in the architecture and
to allow recognition to depend on the particular HLA context. It is not to
assume that adding a head automatically removes the training-data bias: an
unconstrained triple head can still learn a conjunction-like rule when no
counterexamples are available. Verified \texttt{110} examples are therefore
needed to test whether this representational freedom is used to distinguish
positive pairs from a positive specified triple. Section~\ref{sec:method:problem}
separates this modelling rationale from its empirical validation.

Unifying the three tasks is not itself new. UniPMT already jointly predicts
peptide--MHC, peptide--TCR and peptide--MHC--TCR binding using a heterogeneous
graph \citep{zhao2025unipmt}. Our first contribution is a different unified
formulation: a sequence-based model with three supervised heads and explicit
input boundaries, in which \PH{} uses $(p,h)$, \PT{} uses $(p,t)$ and
\PHT{} uses $(p,h,t)$. UniPMT's \PT{} score is obtained by averaging its
\PHT{} predictions over candidate MHC molecules. In contrast, our directly supervised \PT{} branch
uses only the peptide and TCR sequences and supplies its unpooled residue
features to an explicit \PHT{} branch. The triple head composes these with
\PH{} features within the current sample, without a binding-relation graph.
This is a distinction in representation and supervision design, not a claim to be the first unified
predictor or a claim of superiority without comparison.

The second contribution concerns how such a model can share capacity while
accommodating competing task requirements. UnifyImmun shares a peptide encoder
but uses separate cross-attention modules for \PH{} and \PT{}, trained with a
two-phase progressive schedule \citep{yu2025unifyimmun}. Its described
training objective has two pairwise heads, rather than an independently
supervised triple head. Fusion-pMT instead pretrains a \PH{} model, passes its
unpooled interaction representation into a TCR attention module, and finetunes
the model on \PHT{} \citep{ma2025fusionpmt}. Neither design directly specifies
how to balance all three objectives when they remain active together. Extending
the former with a triple head exposes the shared encoder to three objectives
but leaves interaction-level sharing unspecified. Extending the latter with
simultaneous pair supervision makes the reused \PH{} representation serve
both HLA binding and TCR recognition. These are potential sources of training
interference, not reported failures of either published method.

The tension is concrete. A \PH{} predictor needs features that work across
TCRs because it never receives a TCR; a \PT{} predictor must work without a
specified HLA; a \PHT{} predictor must distinguish the particular combination.
A representation useful for a pair task can therefore omit information useful
for the triple task. Full sharing can make those objectives compete for the
same transformation, whereas completely separate interaction modules cannot
reuse transformations learned from the other relations. We address this
trade-off with two globally shared experts and two private experts in each
relation branch, combined by a learned relation-specific gate. Shared experts
provide reusable transformations; private experts provide alternatives for
relation-specific requirements; and the gate adjusts the mixture to the
sample. We adopt this topology from multi-task learning \citep{tang2020ple}
and adapt it to the peptide residue axis, so that branches with different
inputs can use the same expert pool and retain their outputs for subsequent
triple interaction. This provides a mechanism for selective sharing, not a
guarantee that all gradient conflicts disappear.

Our transfer contribution starts from Fusion-pMT's use of a \PH{}
representation learned under pairwise supervision as input to the triple
interaction \citep{ma2025fusionpmt}. This captures the presentation side of
recognition, but a triple decision depends on both peptide--HLA binding and
TCR recognition of the presented complex. Fusion-pMT learns the latter from
the triple objective; it does not additionally transfer interaction features
learned by an independently supervised \PT{} task. We add that complementary
path. An HLA-conditioned gate injects the unpooled \PT{} representation at
each peptide position, so the triple branch receives both presentation
features and independently learned recognition features, with their use
adjusted to the current HLA context.

The connection also benefits the source task through training. Because the
transfer is differentiable, the \PHT{} classification loss propagates back
to the \PT{} interaction representation, making it useful for both pairwise
association and recognition in a specified HLA context. This feedback updates
the \PT{} feature extractor, not its separate classification head through the
transfer path, and \PT{} inference still requires only peptide and TCR inputs.
Thus the mechanism combines \PT{}$\rightarrow$\PHT{} forward feature transfer
with \PHT{}$\rightarrow$\PT{} backward supervision. Our ablation experiments
show improved prediction for both tasks when the gated connection is added;
isolating the contribution of backward feedback requires the gradient-blocking
control in Section~\ref{sec:results:ablation}.

Finally, sharing and transferring features do not specify how the resulting
scores should relate. Separate binary losses train each head against its own
label, but never directly measure a gap between a pair score and the triple
score for the same sample. A positive triple can consequently receive a high
score while an associated positive pair remains under-confident. Simply
bounding the triple score by both pair scores can pass that pairwise error
into the triple decision. Moreover, the heads have different conditioning
inputs, so the label implication on annotated triples does not automatically
justify a pointwise ordering of their predicted probabilities.

We therefore introduce label-aware directional consistency at the output
level. On a positive triple, the penalty raises a lagging positive pair score
using the triple score as a detached reference. When both a pair and the
triple are labelled negative, it instead lowers an excessive triple score
using the pair score as the reference. Each direct correction agrees with
the updated output's binary label; a tolerance margin and finite loss weight
make the coupling a soft training preference rather than an inference-time
constraint. MoE allocates interaction capacity, gated transfer exchanges
features and their supervision, and consistency regulates score relationships.
Together these designs form \method{}
(Figures~\ref{fig:teaser} and~\ref{fig:arch}).

\TODO{Complete the experiment summary with the confirmed transfer-ablation
metrics, splits and seed information; add other findings only when supported
by the corresponding results in Section~\ref{sec:results}.}

Our contributions are as follows.
\begin{enumerate}[1.]
\item \textbf{Unified sequence-based prediction of three relations.} We
jointly supervise \PH{}, \PT{} and \PHT{} heads with their respective input
pairs or triple, and preserve a separate HLA-conditioned triple decision.
Unlike UniPMT's graph-based unification and MHC-averaged \PT{} score, our
\PT{} branch is evaluated directly from its two sequences and its unpooled
features enter the triple head. The separate triple decision avoids imposing
pairwise sufficiency as an architectural constraint and can express
HLA-dependent outcomes beyond conjunction. Whether that capability is learned
remains to be tested with the currently missing \texttt{110} supervision.
\item \textbf{Relation-aware expert sharing for joint pair and triple
learning.} We place an established shared--private MoE topology inside each
interaction block, with sample-level routing over residue-level experts.
Relation-conditioned gates let branches with different input compositions
choose reusable and branch-specific transformations, and the transformed pair
representations continue into the triple branch. This design targets training
interference while retaining transfer; its value is to be evaluated against matched
sharing and placement controls.
\item \textbf{HLA-conditioned gated transfer with reciprocal supervision.}
Fusion-pMT transfers independently learned \PH{} features into \PHT{}; we
add independently supervised \PT{} interaction features as complementary
recognition evidence. For \PHT{}, a per-residue gate adapts these features to
the specified HLA and triple state. For \PT{}, the differentiable connection
returns triple-loss gradients to its feature extractor while preserving its
pairwise input boundary. Forward transfer and backward feedback jointly train
the reused representation; ablation experiments show that the connection
improves both \PT{} and \PHT{} prediction.
\item \textbf{Label-aware directional consistency of the output scores.}
To address score disagreement not explicitly controlled by MoE or feature
transfer, four label-conditioned penalties select both the applicable
pair--triple relation and the direction of correction. Positive triples guide
lagging pair scores upward; negative pairs guide excessive negative-triple
scores downward. Stop-gradient holds the reference operand fixed for that
penalty, while a margin and finite weight retain flexibility. Unlike a hard
probability hierarchy, this objective neither caps every triple prediction by
the pair scores nor enforces a universal ordering between differently
conditioned heads.
\item \TODO{Complete the empirical contribution with numerical evidence for
the confirmed PT/PHT transfer gains and the remaining evaluated settings.}
\end{enumerate}

\begin{figure}[!t]
\centering
\figph{5.2cm}{Figure~1 placeholder (teaser). Compare task coverage and
information flow: UnifyImmun uses two pairwise branches over a shared peptide
encoder; Fusion-pMT transfers an unpooled PH representation into PHT after
pair pretraining; UniPMT jointly predicts three relations through a graph.
UniReMoE jointly supervises three sequence-based heads, shares interaction
experts selectively and transfers unpooled PT features under an HLA-conditioned
gate. Mark missing class 110 as an evaluation gap.}
\caption{Task unification and selective sharing address distinct questions:
which relations are predicted, and how their interaction representations are
shared. The proposed model retains a separate triple head and combines shared
and branch-private experts with residue-level transfer.}
\label{fig:teaser}
\figalttext[Task coverage and information flow across the compared models]{\TODO{alt text for the completed teaser figure}}%
\end{figure}

%%%%%%%%%%%%%%%%%%%%%%%%%%%%%%%%%%%%%%%%%%%%%%%%%%%%%%%%%%%%%%%%%%%%%%%%%%%%%%
\section{Related work}\label{sec:related}
%%%%%%%%%%%%%%%%%%%%%%%%%%%%%%%%%%%%%%%%%%%%%%%%%%%%%%%%%%%%%%%%%%%%%%%%%%%%%%

\subsection{Peptide--HLA presentation prediction}\label{sec:related:ph}

Presentation prediction is mature and sets the accuracy expectation for the
\PH{} task. Allele-conditioned neural predictors trained on pooled binding
affinity and mass-spectrometry eluted-ligand data generalise across thousands
of alleles by encoding the allele as a pseudo-sequence rather than as a
categorical label \citep{reynisson2020netmhcpan, odonnell2020mhcflurry}, and
transformer variants have since improved peptide-length handling and
pan-allele transfer \citep{chu2022transphla}. We adopt the same
pseudo-sequence view of the allele, and we treat this literature as the
reference point for the \PH{} column of our tables rather than as a component
to be improved.

\subsection{Peptide--TCR recognition prediction}\label{sec:related:pt}

Recognition prediction is the data-limited half of the problem, and this
scarcity is what motivates sharing across tasks. Methods differ mainly in how
they encode the CDR3 and how they score compatibility: distance- and
clustering-based approaches group TCRs by sequence similarity to a known
binder \citep{glanville2017identifying, dash2017quantifiable}, while supervised
models learn a joint peptide--CDR3 scoring function with recurrent, attention or
transfer-learning encoders \citep{springer2020ergo, weber2021titan}. Reported accuracy drops sharply for peptides and TCRs unseen
in training, which is why we evaluate every task under explicit unseen-entity
splits rather than on random partitions.

\subsection{Unified models over both binding events}\label{sec:related:unified}

UniPMT is the closest precedent in task coverage: it jointly learns \PH{},
\PT{} and \PHT{} prediction with a heterogeneous graph
\citep{zhao2025unipmt}. GraphSAGE produces entity representations; a dedicated
\PHT{} predictor combines peptide--MHC and MHC--TCR interaction vectors, and
its training includes triple supervision. The \PT{} score is the arithmetic
mean of \PHT{} scores over candidate MHC molecules (their Eq.~11), and the
three objectives are jointly optimised (their Eq.~12). Thus neither three-task
coverage nor a supervised triple predictor is new to our work. Our distinction
is a direct sequence-based branch for each relation: \PT{} requires no
averaging over candidate HLA molecules and retains an independently supervised,
unpooled peptide--TCR representation. That representation can then be used
inside the HLA-conditioned triple interaction. This comparison concerns the
computation and information flow; it does not imply that UniPMT cannot handle
unseen entities, which it evaluates explicitly.

UnifyImmun provides a complementary precedent for sharing across the pair
tasks. It encodes the three entity types separately and uses two cross-attention
branches over a shared peptide encoder, with a two-phase progressive training
schedule \citep{yu2025unifyimmun}. Although the paper describes applications
to HLA--peptide--TCR prediction, its specified architecture and loss train
\PH{} and \PT{} heads, without a separately supervised triple head. Adding
such a head would require deciding which interaction features to share and
how to preserve pair-task performance under triple supervision. Our MoE
addresses that allocation question. It is not justified by claiming that
UnifyImmun shares one interaction transformation: its two interaction modules
are separate. Nor are UniPMT's outputs independent, since its \PT{} score is
derived from \PHT{} scores. Our label-aware directional objective is a
specific additional coupling mechanism, not the first dependency between
outputs in a unified model.

\subsection{Hierarchical transfer from a pair task to the
triple}\label{sec:related:transfer}

Triple predictors already use the HLA context explicitly. For example,
pMTnet accepts a peptide, an HLA allele and a TCR CDR3$\beta$ sequence to
predict TCR recognition of a peptide--MHC complex \citep{lu2021pmtnet}.
A related line transfers a trained \PH{} model into a triple predictor and
provides the direct precedent for our representation-transfer design. Fusion-pMT pretrains on \PH{}, then discards the pooling and
classification head and reuses the \emph{unpooled} peptide--HLA interaction
matrix as the query of a second cross-attention module against the receptor,
finetuning all parameters on the triple task \citep{ma2025fusionpmt}. The
paper by \citet{niu2025attention} transfers its pretrained peptide--HLA encoder
into a recognition model and adds a margin-based differential objective between
a binding and a non-binding receptor for the same peptide--HLA complex. StriMap
goes further on the molecular side: it pretrains a \PH{} predictor, then
\emph{freezes} it and reuses it as a feature aggregator, mixing its pretrained
representation with a task-specific one through an input-dependent gate on the
HLA side and the peptide side separately, and models residue--residue
interaction with bilinear attention over paired receptor $\alpha$/$\beta$
chains and ESMFold-predicted structures \citep{cao2026strimap}.

We adopt the central insight of this line --- that the unpooled pair
representation, not a pooled score, is what the triple task should consume ---
and we do not claim it. Our increments over it are three, and each is narrow.
First, all three of these methods transfer in one direction only, from \PH{} to
the triple, and none of them has an independently supervised peptide--TCR
branch; we add one, so that the triple task receives recognition evidence that
was itself trained against \PT{} labels rather than inferred from the triple
objective alone. Second, Fusion-pMT and \citet{niu2025attention} replace the \PH{} training
objective with triple training in a second stage, whereas StriMap freezes its
\PH{} predictor. Fusion-pMT's reused parameters do receive triple gradients;
what is absent is a continuing \PH{} objective that preserves pair-task
performance during that adaptation. We keep all three objectives active and
both transfer paths differentiable, so pair and triple supervision act on the
reused modules together. In particular, the added \PT{} path supplies
recognition features to \PHT{} and returns its classification gradients to
the \PT{} extractor. Our ablation shows improvements in both tasks from adding
the gated connection; this does not by itself isolate the roles of gating,
forward reuse and backward feedback. Third, StriMap's gate mixes two representations \emph{of the same
entity} --- a frozen pretrained one against a task-specific one --- whereas our
gate controls how much of a \emph{different relation's} representation enters
the triple path at each peptide residue.

Two limits of that comparison must be stated plainly. StriMap consumes paired
$\alpha$/$\beta$ chains and predicted three-dimensional structure, while we use
a single CDR3$\beta$ loop and no structural input, so we cannot claim more
complete modelling of the molecular interaction; our claim is confined to how
supervision is used and how representations are transferred. And the
differential objective of \citet{niu2025attention} compares two samples within
one task, whereas our consistency terms compare three outputs of one sample
across tasks; the two are complementary, and ours is not a substitute for
learning from real triple negatives.

Table~\ref{tab:positioning} summarises where each mechanism we build on comes
from and what we change about it. We state the inherited mechanism first in
every row, distinguishing task coverage from the mechanisms used to share
representations and couple predictions.

\begin{table*}[!t]
\caption{Mechanisms this work builds on, the limitation of each as it stands in
the cited work, and the specific change we make. Rows are grouped by which of
our four methodological contributions they bear on. No row claims a mechanism as ours unless
the final column says so.\label{tab:positioning}}
\fontsize{7.5pt}{9pt}\selectfont
\tabcolsep=0pt%
\begin{tabular*}{\textwidth}{@{\extracolsep{\fill}}p{0.16\textwidth}p{0.27\textwidth}p{0.25\textwidth}p{0.27\textwidth}@{\extracolsep{\fill}}}
\toprule
Prior work & Mechanism we take from it & Limitation for our setting
& What we change\\
\midrule
\multicolumn{4}{@{}l}{\textit{Bearing on contribution 1: unified three-relation
prediction}}\\
UniPMT \citep{zhao2025unipmt}
& Joint \PH{}, \PT{} and \PHT{} prediction using a heterogeneous graph.
& The \PT{} score averages \PHT{} scores over candidate MHC molecules;
there is no direct residue-level \PT{} branch to transfer.
& Three direct sequence branches; independently supervised \PT{} residue
features feed the HLA-conditioned triple branch.\\
UnifyImmun \citep{yu2025unifyimmun}
& A shared peptide encoder with separate \PH{} and \PT{} interaction heads.
& Its specified training objective has no separately supervised triple head.
& A jointly supervised \PHT{} head, with selective interaction sharing.\\
\midrule
\multicolumn{4}{@{}l}{\textit{Bearing on contribution 2: shared--private
interaction experts}}\\
MMoE \citep{ma2018mmoe}
& One shared expert pool with a per-task softmax gate; per-task gate
distributions reported as a diagnostic.
& No dedicated private path for relation-specific transformations;
experts act on a pooled vector.
& We keep the per-task gate and the diagnostic, and add private experts.\\
CGC / PLE \citep{tang2020ple}
& Shared expert module plus task-exclusive modules, fused by a dense per-task
gate. \textbf{This is the topology we use.}
& Posited as a general remedy for competing tasks; where its sample space is
heterogeneous, the tasks still read one common input representation and differ
only in label coverage; experts are MLPs on pooled vectors feeding a task
tower; gate reads raw input features.
& We adapt the decomposition to relations with unequal inputs: position-wise
experts on the peptide axis, output continuing into the triple interaction,
gate conditioned on a relation summary.\\
NeoMUST \citep{ma2024neomust}
& The same gate-control topology applied in immunopeptidomics, with
uncertainty task weighting and conflict-averse gradients.
& Both tasks are peptide--MHC over one input pair, so sharing across differing
input scopes is untested; no triple relation.
& Three relations with different inputs. We use no gradient surgery and claim
no advantage in handling gradient conflict.\\
MoE-TCR \citep{qu2026moetcr}
& Expert mixtures for receptor--epitope prediction.
& Experts are whole predictors under a sparse top-$k$ gate routed on the
epitope; nothing is shared at the representation level.
& Experts inside the interaction block, mixed densely, with a shared pool
across relations.\\
\midrule
\multicolumn{4}{@{}l}{\textit{Bearing on contribution 3: residue-level gated
transfer}}\\
Fusion-pMT \citep{ma2025fusionpmt}
& The unpooled peptide--HLA interaction matrix, not a pooled score, is what
the triple task should consume. \textbf{This insight is theirs.}
& One-directional \PH{}$\rightarrow$\PHT{} transfer with no peptide--TCR
branch; triple training replaces pair training in a second stage.
& An HLA-gated \PT{} feature path supplies independently supervised recognition
features to \PHT{} and returns triple gradients to the \PT{} extractor;
pair and triple objectives remain active.\\
\citet{niu2025attention}
& Transfer of a pretrained peptide--HLA encoder into a recognition model.
& Its differential objective compares two samples within one task, and it has
no peptide--TCR branch.
& Our coupling is across the three outputs of one sample; the two are
complementary and we do not replace theirs.\\
StriMap \citep{cao2026strimap}
& Gated fusion of a pretrained pHLA representation with a task representation;
residue-level interaction modelling.
& The pHLA predictor is frozen, so information flows one way and the pair task
cannot benefit from triple supervision.
& Binary interaction modules stay trainable under triple gradients. Our gate
mixes a \emph{different relation's} representation rather than two
representations of the same entity.\\
\midrule
\multicolumn{4}{@{}l}{\textit{Bearing on contribution 4: directional
consistency}}\\
HMCN \citep{wehrmann2018hmcn}
& The squared-hinge violation penalty. \textbf{This form is theirs.}
& Symmetric in gradient: both operands move when it fires, and the label never
enters the penalty.
& Labels select the reference; stop-gradient blocks its direct derivative in
the penalty, so only the label-supported score receives a direct correction.\\
C-HMCNN \citep{giunchiglia2020chmcnn}
& Label-aware gradient routing: the label masks candidates inside a maximum.
& Consistency is architectural, and the direction is fixed by a taxonomy edge,
one direction per edge for every sample.
& Soft penalty with no hard projection, stop-gradient instead of candidate
masking, and a direction chosen per sample from the observed label pattern.\\
\botrule
\end{tabular*}
\begin{tablenotes}%
\item StriMap additionally uses paired receptor $\alpha$/$\beta$ chains and
predicted structures, which we do not; we therefore make no claim of more
complete molecular modelling.
\end{tablenotes}
\end{table*}

\subsection{Protein language models as frozen feature
extractors}\label{sec:related:plm}

Protein language models supply the residue-level representations our
interaction modules operate on, and keeping them frozen is a deliberate
design choice. Masked-language-model pretraining on large sequence corpora
yields per-residue embeddings that transfer to structure and function tasks
\citep{rives2021esm, lin2023esmfold}, and repertoire-pretrained variants
provide the same for TCR sequences \citep{wu2024tcrbert}. Fine-tuning them on
immunopeptidomics-scale data is expensive and risks overfitting the scarce
\PT{} labels, so a common alternative is to freeze the encoder and learn a
small adapter on top \citep{houlsby2019parameterefficient}. We follow that
route and go one step further by precomputing the residue embeddings once into
a fixed-width cache, so the language models are absent from the training graph
entirely.

\subsection{Mixture-of-experts and multi-task
fusion}\label{sec:related:moe}

Mixture-of-experts layers and multi-task gating give us the vocabulary for our
second contribution, but our use of them differs from both traditions in scale
and in purpose. In language modelling, expert layers replace feed-forward
sublayers and are made sparse with top-$k$ routing and capacity limits so that
parameter count can grow without a matching increase in computation
\citep{shazeer2017outrageously, fedus2022switch}. In multi-task learning,
gated expert pools let tasks select different mixtures of a common expert set
and so reduce interference between competing objectives. MMoE shares one
expert pool across all tasks and gives each task its own softmax gate over it,
with no task-private experts \citep{ma2018mmoe}. Customized Gate Control, the
building block of PLE, adds exactly what MMoE lacks: a shared expert module
whose parameters every task updates, plus a task-exclusive module per task,
fused by a dense per-task gate, introduced to counter the ``seesaw
phenomenon'' in which gains on one task come at another's expense
\citep{tang2020ple}. \textbf{The sharing topology we use is theirs.} A pool of
shared experts beside branch-private experts, mixed by a dense per-task softmax
gate, is CGC's design, and the number of experts in each module is a tuned
hyper-parameter in that scheme, so our two-plus-two configuration is a setting
within it rather than a departure from it. Our question is a different one. CGC
posits the shared/specific decomposition as a general remedy for competing
tasks; we ask whether peptide binding across three relations actually has that
structure, where to place the decomposition so relations
with unequal input scopes can use it, and whether the resulting allocation can
be measured and tested against matched sharing controls. Those are questions
about this problem, not about the mechanism, and
Section~\ref{sec:method:moe} answers them. What we change is the carrier of that
topology and what the gate reads. In MMoE and CGC/PLE every task consumes the
same input representation for a given sample, the experts are multilayer
perceptrons over a flat pooled vector, and the mixture output goes to a task
tower or to another extraction layer of the same kind; the gate is a linear
function of that shared input, or in PLE's upper levels of the task's own
previous fusion output. Our three branches instead read different entity
compositions, so the shared experts are updated by branches that do not see the
same inputs; the experts are position-wise and their output keeps the peptide
residue axis, flowing into further interaction computation before any pooling;
and the gate is conditioned on the current relation state and the available
entity summaries, rather than only on raw input features.

The first of those three needs one clarification, because PLE is also evaluated
on tasks with a heterogeneous sample space, in which a user must click an
impression before a share or a conversion can be observed
\citep{tang2020ple}. That heterogeneity sits in the \emph{label} space and is
nested: every task reads the same user--item features, and the tasks differ in
which samples carry a label at all. Ours sits in the \emph{input} space and is
not nested --- the \PH{} branch never receives a receptor and the \PT{} branch
never receives an allele --- so a shared expert is updated by gradients
computed from representations built over different entities, not from the same
representation under different label coverage. Whether a shared expert pool
stays useful under that condition does not follow from the recommendation
results, and it is one of the things Section~\ref{sec:results:ablation} tests.
We should be careful with the second of the three: placing a mixture in the feed-forward sublayer of an
attention block is standard in the sparse-expert transformer line cited above,
so what is particular here is its combination with a shared/private split and
a relation-conditioned gate, not per-position experts as such. Our setting also
takes the multi-task motivation and not the scaling one: the experts are small,
the mixture is dense, and the purpose is to separate relation-general from
relation-specific transformations rather than to add capacity cheaply.
In our own domain the closest
use of this machinery is NeoMUST, which places a customized gate control layer
--- one shared expert and two task-exclusive experts --- over two peptide--MHC
tasks, presentation classification and affinity regression, and adds
uncertainty-based task weighting and conflict-averse gradient descent
\citep{ma2024neomust}. Its two tasks nonetheless read the same input pair, so
its machinery is never asked to share capacity between relations with
\emph{different} input scopes, which is the harder case here; we also apply no
gradient surgery, and therefore make no claim to handle gradient conflict
better than it does.

Taken together, these two lines locate what is open. Expert mixtures have been
used in this domain at the level of whole predictors routed by epitope
\citep{qu2026moetcr} and over two peptide--MHC objectives on a shared input
pair \citep{ma2024neomust}, and the shared-plus-private topology itself is
settled \citep{tang2020ple}. Our contribution is its adaptation to
residue-level interactions that serve presentation, recognition and the triple
relation together, despite their different input compositions. The claim
concerns \emph{application and placement}, rather than a new expert topology. We also inherit the diagnostic rather than
invent it: MMoE reports the softmax gate distribution of each task over the
experts and reads off from it how the model balances shared against
non-shared parameters, and PLE reports the mean and standard deviation of every
gate's expert weights over the test data, concluding that its shared experts
exert the larger influence \citep{ma2018mmoe, tang2020ple}. Inspecting per-task
gate mass is therefore established practice. What we add is only its reduction
to one reported scalar per branch, the shared fraction of
Section~\ref{sec:method:moe}, tracked over training and across evaluation sets.

Closest to our placement,
expert layers have been applied to TCR specificity prediction
\citep{qu2026moetcr}, but there the experts are whole pan-specific predictors
selected by a sparse top-$k$ gate conditioned on the epitope. We instead place
the mixture inside each binding-interaction block, on residue-resolved
representations, combine all experts densely, and share part of the expert pool
across relations.

\subsection{Consistency and logical constraints in multi-task
prediction}\label{sec:related:consistency}

Our fourth contribution belongs to a line of work that injects known relations
between outputs into the loss rather than into the architecture, and two
hierarchical classifiers in that line are close enough that we must state the
difference precisely. Constraint- and logic-based regularisers penalise output
combinations that violate a stated rule by relaxing the rule into a
differentiable surrogate \citep{xu2018semantic, hu2016harnessing}, and
hierarchical classification enforces ordering between a label and its ancestors
\citep{deng2014large}. HMCN supplies the functional form we use: a squared
hinge $\lambda\max\{0, \hat y_{\mathrm{child}} - \hat y_{\mathrm{parent}}\}^2$
penalising a child score that exceeds its parent's \citep{wehrmann2018hmcn}. We
claim no credit for that form. Its penalty is, however, symmetric in gradient:
when it fires, the child is pushed down and the parent pushed up by equal
amounts, and the ground-truth label enters only through the classification
losses, never through the penalty itself. C-HMCNN is closer still, because it
does make the coupling label-aware: its loss multiplies predictions by their
binary labels inside a maximum, so the label chooses which outputs are eligible
to receive gradient and the maximum selects the recipient among them, while a
max-constraint module enforces hierarchy consistency architecturally
\citep{giunchiglia2020chmcnn}.

Our rule differs from both on three specific points, and on nothing else.
First, in form: we add soft squared hinges to the objective and enforce no hard
constraint, where C-HMCNN's consistency is architectural and its label
awareness sits inside a likelihood. Second, in the mechanism of gradient
control: C-HMCNN masks a candidate set by multiplying by the label, whereas we
apply an explicit stop-gradient to the operand the label designates as the
reference, which keeps that operand in the forward value while freezing it as a
target. Third, in the choice of correction target: for a given pair--triple
relation, our observed label pattern selects whether the pair score is raised
or the triple score is lowered, and when the coupling is inactive. This changes
the gradient recipient, not the direction of a biological implication. Our
contribution is this label-dependent application of soft directional coupling
to the three relation heads, rather than a new hinge function or a new claim
that labels can guide optimisation.

%%%%%%%%%%%%%%%%%%%%%%%%%%%%%%%%%%%%%%%%%%%%%%%%%%%%%%%%%%%%%%%%%%%%%%%%%%%%%%
\section{Materials and methods}\label{sec:method}
%%%%%%%%%%%%%%%%%%%%%%%%%%%%%%%%%%%%%%%%%%%%%%%%%%%%%%%%%%%%%%%%%%%%%%%%%%%%%%

\subsection{Overview}\label{sec:method:overview}

\method{} predicts the three binding relations with one network over frozen
protein-language-model residue embeddings, and it is organised so that every
representation keeps the peptide residue axis until the final pooling step.
Figure~\ref{fig:arch} shows the architecture. Three trainable sequence encoders
map cached residue embeddings of the peptide, the HLA pseudo-sequence and the
CDR3$\beta$ loop into a common $\dmodel$-dimensional space
(Section~\ref{sec:method:plm}). Three interaction blocks, one per relation,
then let the peptide attend to its partner and pass the result through a gated
mixture of shared and branch-private experts
(Section~\ref{sec:method:moe}). The triple block additionally receives the
unpooled \PT{} representation through a per-residue, HLA-conditioned gate;
this differentiable connection also carries triple-loss gradients back to the
\PT{} feature extractor (Section~\ref{sec:method:fusion}). Three attention-pooling heads produce the
three logits, which are coupled during training by four label-aware
directional hinges (Section~\ref{sec:method:consistency}) within the overall
objective of Section~\ref{sec:method:loss}. Sections
\ref{sec:method:data}--\ref{sec:method:impl} give the data protocol and the
implementation settings.

\begin{table*}[!t]
\caption{Distinct roles of the three coupling mechanisms. These are structural
roles, not claims of independently established performance gains.
\label{tab:mechanism-roles}}
\tabcolsep=0pt%
\begin{tabular*}{\textwidth}{@{\extracolsep{\fill}}p{0.16\textwidth}p{0.25\textwidth}p{0.27\textwidth}p{0.27\textwidth}@{\extracolsep{\fill}}}
\toprule
Mechanism & Problem addressed & Operation & What it does not enforce\\
\midrule
Shared--private MoE
& Different relations need reusable and specialised interaction capacity.
& Relation-specific routers mix shared and private residue-level transformations.
& No direct rule for pair--triple score gaps; no gradient-conflict guarantee.\\
PT-to-PHT gated transfer
& The triple decision can reuse independently supervised recognition features.
& HLA-conditioned feature injection with backward feedback to the PT extractor.
& No explicit score ordering or choice of which output to correct.\\
Label-aware consistency
& Separate heads can retain score gaps despite shared representations.
& Labels select an eligible coupling and correction target; a detached
reference anchors a soft penalty.
& No new feature-transfer path, exact hierarchy or guaranteed calibration.\\
\botrule
\end{tabular*}
\end{table*}

\begin{figure*}[!t]
\centering
\figph[\textwidth]{7.8cm}{Figure~2 placeholder (architecture). Panel (a):
frozen ESM-C and TCR-BERT produce cached per-residue embeddings; three
trainable encoders project them to $\dmodel=256$. Panel (b): the \PH{} and
\PT{} interaction blocks (peptide as query, no feed-forward sublayer) each feed
a \texttt{BranchMoE} of two shared plus two branch-private experts under a
temperature-annealed softmax gate. Panel (c): the triple block takes the
post-expert \PH{} representation as query and the TCR as key/value, then
receives $\beta_i R(Z^{\PTm}_i)$ added at every peptide residue, where $R$ is
an affine projection after layer normalisation, with
$\beta_i$ gated on the triple state, the \PT{} state and the pooled HLA
summary. Show the forward PT-to-PHT feature arrow and a distinct backward
PHT-loss-to-PT-extractor gradient arrow; do not draw a symmetric forward
connection or a transfer-gradient arrow into the PT classifier.
Panel (d): three attention-pooling heads and the four directional
hinges linking their outputs. Mark clearly which gate is sample-level
($g\in\Delta^3$) and which is per-residue ($\beta\in[0,1]^{L_P}$).}
\caption{\TODO{architecture caption: name every module, state tensor shapes,
and state that the two gates act at different granularities.}}\label{fig:arch}
\figalttext[Architecture of the proposed model]{\TODO{alt text for the architecture figure}}%
\end{figure*}

\subsection{Problem formulation}\label{sec:method:problem}

We learn three binary prediction functions with one parameterised model. Let
$p$ be a peptide of length $L_P\in[8,15]$, $h$ an HLA class-I pseudo-sequence of
up to 34 residues and $t$ a CDR3$\beta$ loop of up to 34 residues. The scores
$\hat y_{\PHm}(p,h)$, $\hat y_{\PTm}(p,t)$ and
$\hat y_{\PHTm}(p,h,t)$ respectively predict peptide--HLA binding/presentation,
peptide--TCR compatibility without an explicit HLA input, and recognition of
the specified peptide--HLA--TCR combination. These are different conditioning
sets, not three interchangeable estimates of the same probability. Separate
pair predictions require no placeholder for the missing entity. The model
uses only CDR3$\beta$ for the receptor and does not represent the full TCR
contact surface.

\begin{definition}[Triple label class]\label{def:class}
For a triple row with all three labels observed, its class is the three-digit
code $y_{\PHm}y_{\PTm}y_{\PHTm}$.
\end{definition}

The complete triple-labelled rows have support
$\mathcal{S}_{\mathrm{obs}}=\{000,010,100,111\}$. Here \texttt{100} has
positive \PH{} evidence but negative \PT{} and \PHT{} labels; \texttt{010}
has a positive \PT{} association but negative \PH{} and \PHT{} labels for
the specified HLA. Neither description implies recognition of a free peptide.
Under this corpus's annotation convention, a positive triple also has positive
pair labels, excluding \texttt{001}, \texttt{011} and \texttt{101}. The
additional absence of \texttt{110} implies
\begin{equation}
y_{\PHTm}=y_{\PHm}\land y_{\PTm},
\qquad (y_{\PHm},y_{\PTm},y_{\PHTm})\in\mathcal{S}_{\mathrm{obs}}.
\label{eq:observed-conjunction}
\end{equation}
This is a property of the observed labels, not a general law of recognition.
In particular, observing \texttt{111} does not establish the reverse
implication from pairwise positivity to recognition under every HLA context.
Verified \texttt{110} cases would distinguish the two prediction questions.

\paragraph{Why retain a separate triple task?}
Equation~\eqref{eq:observed-conjunction} describes the available supervision;
we do not impose it as the model's prediction rule. Instead of defining
$\hat y_{\PHTm}$ through a fixed combination of $\hat y_{\PHm}$ and
$\hat y_{\PTm}$, we learn it from the joint interaction representation of
$(p,h,t)$ with its own supervised objective. The head can therefore retain
information about the specified HLA context that is not captured by two scalar
pair scores. This provides a way to represent triple-specific outcomes without
assuming that positive pair predictions are sufficient.

Representational freedom and learned behaviour must be distinguished. With
only the four observed classes, a separate head can still learn the same
conjunction-like regularity; its presence alone does not remove that bias.
Moreover, the present data reader and loss validator reject \texttt{110}.
Testing whether the head learns exceptions requires verified examples and an
extension of the data/loss contract. We compare
the triple score with $s_{\min}=\min(\hat y_{\PHm},\hat y_{\PTm})$ and
$s_{\mathrm{prod}}=\hat y_{\PHm}\hat y_{\PTm}$. Even though the true labels
satisfy Eq.~\eqref{eq:observed-conjunction}, imperfect pair predictions need
not yield the same ranking or calibration as a learned triple head. A gain
over these controls would show predictive value on the available support,
not establish discrimination of the unobserved \texttt{110} class.

\subsection{Frozen residue representations and sequence
encoders}\label{sec:method:plm}

We keep both protein language models frozen and outside the training graph, so
that all learned capacity sits in the downstream network. Peptide and HLA
residues are embedded with ESM-C (300M parameters, hidden size 960)
\citep{esmc2024} and the CDR3$\beta$ loop with TCR-BERT (hidden size 768)
\citep{wu2024tcrbert}. Both encoders are run once in inference mode over the
whole corpus and their per-residue outputs are written to a sharded,
fixed-width \texttt{float16} HDF5 store indexed by sequence, so training reads
tensors from disk and never instantiates a language model. No language-model
parameter is therefore trainable, or even present, during training.

Three trainable encoders then bring the three entity types into a common space.
For entity $e\in\{P,H,T\}$ with cached features $X_e\in\mathbb{R}^{L_e\times
d_e}$, a linear projection with layer normalisation and dropout maps $d_e$ to
$\dmodel=256$, three pre-norm transformer encoder layers
($8$ heads, feed-forward width $1024$, GELU) model within-entity context, and a
final layer normalisation gives $H_e\in\mathbb{R}^{L_e\times \dmodel}$. We add
no positional encoding, because the cached language-model features already
carry position. The peptide encoder is shared by all three relations, the HLA
encoder by \PH{} and \PHT{}, and the TCR encoder by \PT{} and \PHT{}. Padding
positions are masked throughout, and we write $\bar H_e$ for the mask-aware
mean of $H_e$ over real residues.

\subsection{Relation interaction with shared--private
experts}\label{sec:method:moe}

\paragraph{Motivation: selective sharing under different task requirements.}
The model must reuse information across three related tasks without requiring
them to use identical transformations. The biological motivation is simple:
HLA holds a peptide in its binding groove and influences how it is displayed;
the TCR contacts the exposed peptide together with the HLA surface
\citep{rudolph2006structural}. The peptide is common to the three predictions,
but the features useful for HLA binding, HLA-unspecified TCR association and
recognition of a particular complex need not be identical. This motivates a
shared--private decomposition; it does not identify individual experts with
particular molecular contacts.

The training problem follows from the input boundaries. The \PH{} branch must
remain useful without knowing the TCR, and \PT{} without knowing the HLA,
whereas the triple loss can reward features specific to their combination.
If all interaction capacity is shared, updates useful for one objective may
impair another; if all interaction capacity is separate, useful transformations
cannot be reused across relations. In a joint extension of Fusion-pMT, this
tension would also arise in the pair representation reused by the triple
branch. In a joint extension of UnifyImmun, adding a triple head would leave
its allocation between shared encoding and separate interactions to be
specified. These are reasons to test selective sharing, not evidence that the
published models suffer measured negative transfer.

We use the shared--private topology of CGC/PLE \citep{tang2020ple} at the
interaction level. Its components address three requirements. Shared experts
allow information learned from one relation to support the others. Private
experts give each relation additional transformations that the other branches
are not forced to apply. Relation-conditioned gates adjust their mixture for
the current sample, rather than fixing one amount of sharing for every peptide
and partner. All experts operate on the peptide residue axis, a common
interface available to all three branches despite their different inputs.
Keeping that axis also allows the transformed pair features to enter the
subsequent triple interaction. Sharing after pooling is possible, but would
remove the position-specific features from that transfer path.

This allocation is branch-specific, not gradient isolation. The triple loss
still reaches shared encoders and experts, and also the \PH{}/\PT{} private
experts through the reused pair representations. The MoE provides alternatives
to mandatory sharing and is intended to mitigate interference; it neither
performs gradient surgery nor guarantees the absence of conflict. The sharing
controls and single-task comparisons in Section~\ref{sec:results:ablation}
are needed to establish whether that design helps.

\paragraph{Interaction block.}
Each relation $r\in\{\PHm,\PTm,\PHTm\}$ has one cross-attention block in
which the peptide is always the query, so that the output keeps the peptide
residue axis. With query $Q$, key/value $K$ and the key padding mask applied
inside the attention,
\begin{equation}
U_r \;=\; Q \;+\; \mathrm{Drop}\!\left(\MHA\!\left(\LN(Q),\,\LN(K),\,\LN(K)\right)\right),
\label{eq:attn}
\end{equation}
where $\LN$ is layer normalisation, the residual is taken on the
un-normalised query, and padding positions of the output are set to zero. For
$r=\PHm$ we set $Q=H_P,\,K=H_H$ and for $r=\PTm$ we set $Q=H_P,\,K=H_T$; the
triple case is given in Section~\ref{sec:method:fusion}. Unlike a standard
transformer block, Eq.~\eqref{eq:attn} has no position-wise feed-forward
sublayer: that sublayer is exactly what the expert mixture replaces.

\paragraph{Shared--private expert mixture.}
Each block owns a mixture of four position-wise experts, two of which are
shared with the other relations. An expert is a two-layer feed-forward network
$F(z)=W_2\,\mathrm{GELU}(W_1 z)$ with inner width $512$. The model holds one
global pool $\{F^{\mathrm{S}}_1,F^{\mathrm{S}}_2\}$ registered once and used by
all three relations, and each relation $r$ additionally owns
$\{F^{r}_1,F^{r}_2\}$. A relation-specific router produces one mixture weight
vector per sample from a context vector $c_r$,
\begin{equation}
g_r \;=\; \softmax\!\left(\frac{R_r(c_r)}{\tau}\right)\in\Delta^{3},
\qquad
R_r:\mathbb{R}^{m_r}\!\rightarrow\!\mathbb{R}^{4},
\label{eq:gate}
\end{equation}
with $R_r$ a two-layer GELU network of inner width $128$ and $\tau$ a routing
temperature. The block output applies the mixture at every peptide position
$i$,
\begin{equation}
Z_{r,i} \;=\; U_{r,i} \;+\; \mathrm{Drop}\!\left(\sum_{j=1}^{4}
g_{r,j}\,F_{r,j}\!\left(\LN(U_{r,i})\right)\right),
\label{eq:moe}
\end{equation}
where $(F_{r,1},\dots,F_{r,4}) = (F^{\mathrm{S}}_1,F^{\mathrm{S}}_2,
F^{r}_1,F^{r}_2)$ and padding positions are zeroed. The mixture is dense: all
four experts are evaluated and combined, with no top-$k$ selection and no
capacity limit. The gate is computed once per sample and shared by all
positions, so routing is sample-level while expert computation is
residue-level; we use that wording throughout and never describe the routing
as residue-level.

\paragraph{Router context and task boundaries.}
The router sees a summary of the relation it serves and nothing it should not.
We build $c_r$ by concatenating mask-aware means,
\begin{equation}
\begin{aligned}
c_{\PHm} &= [\,\bar U_{\PHm};\,\bar H_P;\,\bar H_H\,], \\
c_{\PTm} &= [\,\bar U_{\PTm};\,\bar H_P;\,\bar H_T\,], \\
c_{\PHTm} &= [\,\bar{\tilde U}_{\PHTm};\,\bar H_P;\,\bar H_H;\,\bar H_T\,],
\end{aligned}
\label{eq:context}
\end{equation}
giving $m_r=3\dmodel$ for the pair relations and $4\dmodel$ for the triple. The
\PH{} router therefore never reads TCR information and the \PT{} router never
reads HLA information, so each pair prediction depends only on the entities
that relation is defined over. No router reads its own expert output.

\paragraph{Why the placement fits the three tasks.}
The common peptide axis allows the same shared experts to process \PH{},
\PT{} and \PHT{} interaction states without inventing a missing TCR or HLA
input. The \PH{} expert output $Z_{\PHm}$ becomes the query of the triple
attention block, and $Z_{\PTm}$ is transferred before pooling
(Section~\ref{sec:method:fusion}). Expert selection therefore changes the
representations used in downstream interaction, not just the final classifier.
The \PHT{} branch has its own private experts and router to adapt the fused
state to the triple decision. The expert gate chooses \emph{which
transformations to use}; the transfer gate chooses \emph{how much \PT{}
evidence to inject at each residue}. These two decisions serve different roles.

We measure the shared fraction
\begin{equation}
\alpha_r \;=\; g_{r,1}+g_{r,2}
\label{eq:alpha}
\end{equation}
to describe the gate mass assigned to reusable transformations. If learned
allocation is useful, we expect meaningful variation across relations and
samples, and a benefit over capacity-matched fixed or undifferentiated sharing.
Sections~\ref{sec:results:routing} and~\ref{sec:results:ablation} test these
expectations using routing statistics and controlled variants. Gate variation
alone does not demonstrate improved prediction, reduced gradient conflict or
biological specialisation. Inspecting gate weights is an established diagnostic
\citep{ma2018mmoe, tang2020ple}, and $\alpha_r$ is a summary of that diagnostic,
not a biological contribution ratio. The mixture evaluates all four experts,
so we make no sparse-computation efficiency claim.

\subsection{HLA-conditioned gated transfer and reciprocal
supervision}\label{sec:method:fusion}

\paragraph{Motivation: complementing PH transfer with PT evidence.}
Fusion-pMT demonstrates the value of passing interaction features learned by
a supervised \PH{} task into the triple model: the unpooled \PH{}
representation becomes the query for subsequent TCR attention
\citep{ma2025fusionpmt}. This design reflects the dependence of recognition
on HLA presentation. A triple decision also depends on whether the TCR can
recognise the presented complex, so presentation features alone are not the
only reusable source of evidence. Fusion-pMT models TCR recognition under its
triple loss, but has no independently supervised \PT{} branch from which to
transfer recognition features. Our additional \PT{}$\rightarrow$\PHT{} path
fills that specific gap; it does not imply that Fusion-pMT ignores the TCR.

We transfer the unpooled interaction features learned under \PT{} supervision,
alongside the existing \PH{} representation. Because \PT{} has no explicit
HLA input, its features should not be treated as equally informative in every
HLA context. A per-residue gate reads the triple state, the \PT{} state and
an HLA summary to regulate their injection. Joint training keeps this path
differentiable, so it serves two purposes: recognition features support the
\PHT{} decision in the forward pass, and triple supervision refines the
\PT{} feature extractor in the backward pass. This is forward transfer with
backward feedback, not two symmetric forward prediction paths.

\paragraph{Triple path.}
The triple block takes the post-expert \PH{} representation as its query, so
that presentation evidence is carried forward rather than recomputed,
\begin{equation}
U_{\PHTm} \;=\; Z_{\PHm} \;+\;
\mathrm{Drop}\!\left(\MHA\!\left(\LN(Z_{\PHm}),\,\LN(H_T),\,\LN(H_T)\right)\right),
\label{eq:pht-attn}
\end{equation}
which is Eq.~\eqref{eq:attn} with $Q=Z_{\PHm}$ and $K=H_T$.

\paragraph{Gated transfer.}
We then add a projected copy of the \PT{} representation of the same sample at
each peptide residue, scaled by a learned per-residue gate. For every real
peptide position $i$,
\begin{align}
\beta_i &= \sigma\!\left(f\!\left(\left[\LN\!\left(U_{\PHTm,i}\right);\,
\LN\!\left(Z_{\PTm,i}\right);\,\bar H_H\right]\right)\right),
\label{eq:beta}\\
\tilde U_{\PHTm,i} &= U_{\PHTm,i} \;+\; \beta_i\,R\!\left(Z_{\PTm,i}\right),
\label{eq:fuse}
\end{align}
where $f:\mathbb{R}^{3\dmodel}\rightarrow\mathbb{R}$ is a two-layer GELU
network of inner width $128$, and
$R(z)=W_{\PTm}\LN(z)+b_{\PTm}$ is an affine projection with
$W_{\PTm}\in\mathbb{R}^{\dmodel\times\dmodel}$ and
$b_{\PTm}\in\mathbb{R}^{\dmodel}$. Padding positions have $\beta_i=0$.
The final bias of $f$ is initialised to $-2$, favouring weak initial transfer;
$\sigma(-2)\approx0.12$ before the input-dependent contribution to the gate
logit. The fused representation then passes through the triple expert
mixture, $Z_{\PHTm}=\mathrm{MoE}_{\PHTm}(\tilde U_{\PHTm}, c_{\PHTm})$,
following Eqs.~\eqref{eq:gate}--\eqref{eq:context}.

\paragraph{Effect on PHT: recognition features complement presentation features.}
The triple branch receives two sources of pairwise evidence: $Z_{\PHm}$
provides a presentation-conditioned query, and $Z_{\PTm}$ supplies recognition
features learned under their own pairwise labels. The latter need not be
learned solely from triple supervision. Keeping both representations unpooled
lets the triple branch combine them at corresponding peptide positions. The
gate adjusts the contribution of the \PT{} features to the specified HLA and
the current triple state, rather than treating a positive \PT{} association
as sufficient for a positive triple. This supplies complementary recognition
evidence to the \PH{}-derived interaction, and the ablation confirms improved
\PHT{} prediction when the gated connection is included.

\paragraph{Effect on PT: triple supervision refines the interaction representation.}
The same connection makes the \PT{} representation useful to a second
supervised objective. Let $z_i=Z_{\PTm,i}$, $r_i=R(z_i)$ and
$G_i=\partial\mathcal{L}_{\PHTm}^{\mathrm{BCE}}/
\partial\tilde U_{\PHTm,i}$, where
$\mathcal{L}_{\PHTm}^{\mathrm{BCE}}$ denotes the native triple classification
loss. Holding the other inputs of the local fusion module fixed, its gradient
along the transfer edge is
\begin{equation}
\begin{aligned}
\left.\nabla_{z_i}\mathcal{L}_{\PHTm}^{\mathrm{BCE}}\right|_{\mathrm{transfer}}
&= \beta_i J_R(z_i)^{\top}G_i\\
&\quad + (r_i^{\top}G_i)\nabla_{z_i}\beta_i,
\end{aligned}
\label{eq:transfer-feedback}
\end{equation}
where $J_R$ is the Jacobian of $R$. The first term passes through the
transferred feature; the second arises because the gate also reads $z_i$.
$G_i$ includes all downstream uses of the fused state, including its pooled
summary in the triple router context. Thus the extra feedback is not simply
the entire triple gradient multiplied by $\beta_i$.

This path updates the \PT{} interaction block, router and experts, together
with their upstream trainable encoders. It does not pass through the separate
\PT{} pooling or classification head. Those modules remain trained by their
own applicable supervision, while the representation they consume is also
shaped by the triple objective. This additional training requirement can act
as cross-task representation regularisation. Our ablation confirms improved
\PT{} prediction from adding the complete gated connection; attributing that
gain specifically to backward feedback requires the detached-transfer control
in Table~\ref{tab:ablation2}. At inference, \PT{} still receives only $(p,t)$;
training with HLA-conditioned feedback does not give it an HLA input.

\paragraph{Why condition the transfer on HLA?}
Class \texttt{010} illustrates the need for selective use of \PT{} evidence:
the pairwise association is positive, but the specified HLA--peptide pair and
triple are negative. These labels answer different questions and are not
contradictory. An unconditional transfer can make features useful for a
positive \PT{} decision work against a negative \PHT{} decision. The gate
allows the triple branch to reduce their contribution under this HLA context,
while the \PT{} branch retains its own classification objective. This offers
a way to regulate cross-task influence rather than forcing the two decisions
to agree. It does not guarantee gradient alignment: the gate-input term in
Eq.~\eqref{eq:transfer-feedback} and other shared parameter paths remain.

The expert gate $g_r$ is sample-level, whereas the transfer gate $\beta$ is
per-residue. Equations~\eqref{eq:beta}--\eqref{eq:fuse} compose features, not
probabilities, and $\beta$ is not a biological contribution ratio. Forward
feature reuse and backward feedback explain how the connection can benefit
both tasks; the observed joint improvement supports the complete mechanism,
while detached, ungated and explicit-HLA-summary controls separate these
explanations.

\subsection{Label-aware directional consistency}\label{sec:method:consistency}

\paragraph{Why feature sharing still needs output-level coupling.}
MoE and gated transfer connect the representations used by the three heads,
but neither explicitly compares their output scores. Native and auxiliary
binary cross-entropy terms supervise each output separately. They can penalise
an incorrect pair prediction and an incorrect triple prediction, but do not
specify which output should respond to a pair--triple score gap. With finite
data and imperfect fitting, substantial gaps can remain even when the heads
share useful features. We add a relational training signal to address this
remaining problem; it supplements rather than replaces the classification
losses. It adds a training preference, not new labels: a sufficiently well-fitted
model may already produce small score gaps without it.

The labels determine which corrections are appropriate. On \texttt{111}, a
triple score above a pair score motivates raising that positive pair score,
not directly lowering a positive triple score. On \texttt{000}, an excessive
triple score should instead be lowered, not corrected by raising a negative
pair score. On \texttt{010} and \texttt{100}, a positive pair and a negative
triple are legitimate outcomes, not labels that the model should force to
agree. An output coupling should therefore depend on the observed label
pattern and distinguish the reference score from the score to be corrected.
A label-agnostic symmetric penalty does not make these choices: it moves
both operands whenever a gap is penalised, potentially opposing one output's
own label.

\paragraph{Why not impose a hard probability hierarchy?}
A natural alternative enforces
$\hat y_{\PHTm}\le\min(\hat y_{\PHm},\hat y_{\PTm})$ on every prediction.
However, the annotated implication from triple positivity to pair positivity
is not by itself a pointwise law for these three scores. Event inclusion
implies probability ordering under the \emph{same} conditioning information;
our heads instead condition on $(p,h,t)$, $(p,h)$ and $(p,t)$. Even calibrated
predictions under their respective inputs need not obey the proposed ordering
for every triple. Differences in fitting and calibration can further widen
score gaps. We therefore use the hierarchy as a label-conditioned soft prior,
not as a mathematical identity to enforce at inference.

For a concrete hard-projection comparator, define
\begin{equation}
q_{\PHTm}=\min\!\left(\hat y_{\PHTm},\hat y_{\PHm},\hat y_{\PTm}\right).
\label{eq:hard-projection}
\end{equation}
At inference this guarantees the numerical upper bound, but a low pair score
caps the triple score even on a positive triple. It corrects the ordering by
clipping, without using the label to determine whether the low pair score
should instead be raised. If the triple BCE is trained on this minimum, its
direct gradient follows the minimum operand when the minimum is unique; a
pair bottleneck can then cut off that loss's direct gradient to the raw triple
score. Other losses and shared paths can still update the triple branch.
These are trade-offs of this comparator, not proof that every hard hierarchy
performs worse. Our alternative preserves the raw triple prediction and
chooses label-supported corrections during training.

\paragraph{Directional hinges.}
We define four hinge terms, each restricted to the label classes where its
direction is justified, and each treating a label-selected reference operand
as a constant for that penalty. The label selects the correction direction;
it does not certify the reference prediction as accurate. Writing $\hat y_\bullet=\sigma(\ell_\bullet)$ for the
three probabilities, $m$ for a tolerance margin, $\sg[\cdot]$ for the
stop-gradient operator and
\begin{equation}
\phi(a,b)\;=\;\relu\!\left(a-b-m\right)^{2},
\label{eq:hinge}
\end{equation}
the presentation coupling and the recognition coupling are
\begin{align}
\mathcal{L}_{\mathrm{H}} &=
\underbrace{\mathbb{E}_{\mathtt{111}}\!\left[\phi\!\left(\sg[\hat y_{\PHTm}],\,
\hat y_{\PHm}\right)\right]}_{\text{raises }\hat y_{\PHm}}
+\underbrace{\mathbb{E}_{\mathtt{000},\mathtt{010}}\!\left[\phi\!\left(
\hat y_{\PHTm},\,\sg[\hat y_{\PHm}]\right)\right]}_{\text{lowers }\hat y_{\PHTm}},
\label{eq:lh}\\
\mathcal{L}_{\mathrm{T}} &=
\underbrace{\mathbb{E}_{\mathtt{111}}\!\left[\phi\!\left(\sg[\hat y_{\PHTm}],\,
\hat y_{\PTm}\right)\right]}_{\text{raises }\hat y_{\PTm}}
+\underbrace{\mathbb{E}_{\mathtt{000},\mathtt{100}}\!\left[\phi\!\left(
\hat y_{\PHTm},\,\sg[\hat y_{\PTm}]\right)\right]}_{\text{lowers }\hat y_{\PHTm}},
\label{eq:lt}
\end{align}
where each expectation is averaged within its own class group and contributes
zero when the group is empty in a batch. Only rows with all three labels
observed enter these terms.

\paragraph{Label-aligned direct corrections.}
Table~\ref{tab:consistency} specifies the class masks and correction targets.
On \texttt{111}, each active penalty raises a lagging positive pair score and
has no direct derivative through its detached triple reference. On
\texttt{000}, both penalties can lower the negative triple score without a
direct derivative through either pair reference. On \texttt{100}, only the
PT--PHT negative coupling is eligible; on \texttt{010}, only the PH--PHT
negative coupling is eligible. We leave the positive-pair/negative-triple
coupling inactive in these mixed classes. This conservative mask does not
try to equalise legitimately different labels.

Eligibility does not imply a nonzero update: a term fires only when its score
gap exceeds the margin $m$. Let $s_t$ and $s_p$ denote the triple score and
the relevant pair score. For an active positive-class term,
$\phi(\sg[s_t],s_p)$ has a negative derivative with respect to $s_p$, so its
direct correction raises the positive pair score. For an active negative-class
term, $\phi(s_t,\sg[s_p])$ has a positive derivative with respect to $s_t$, so
its direct correction lowers the negative triple score. These directions
agree with binary cross-entropy at the corrected output. The guarantee is
local to the penalty's score derivatives: after an update to shared
parameters, both outputs may change, including the detached reference.

\begin{table}[!t]
\caption{Which directional term is eligible on each label class, and which
output receives its direct gradient when the margin is exceeded. $\uparrow$ raises, $\downarrow$ lowers; an
em-dash means the term is inactive for that class.\label{tab:consistency}}
\begin{tabular*}{\columnwidth}{@{\extracolsep\fill}lccc@{\extracolsep\fill}}
\toprule
Class & $y_{\PHm}y_{\PTm}y_{\PHTm}$
& $\mathcal{L}_{\mathrm{H}}$ acts on & $\mathcal{L}_{\mathrm{T}}$ acts on\\
\midrule
\texttt{111} & 1\,1\,1 & $\hat y_{\PHm}\uparrow$   & $\hat y_{\PTm}\uparrow$\\
\texttt{000} & 0\,0\,0 & $\hat y_{\PHTm}\downarrow$ & $\hat y_{\PHTm}\downarrow$\\
\texttt{100} & 1\,0\,0 & ---                         & $\hat y_{\PHTm}\downarrow$\\
\texttt{010} & 0\,1\,0 & $\hat y_{\PHTm}\downarrow$ & ---\\
\botrule
\end{tabular*}
\begin{tablenotes}%
\item The stop-gradient fixes the probability operand only. It does not freeze
the corresponding branch: the shared encoders, the shared experts and both
transfer paths still receive gradient from the supervised terms.
\end{tablenotes}
\end{table}

\paragraph{Advantages and scope of the soft directional objective.}
The design makes three choices that the hard projection in
Eq.~\eqref{eq:hard-projection} does not. Labels select the applicable relation
and the target of correction, avoiding a direct update against that target's
label. Stop-gradient makes the other operand a reference for this penalty
rather than jointly moving both scores to close the gap. Finally, the margin
ignores small differences, while finite weights allow the classification
objectives to compete with the relational prior. A correct triple prediction
is therefore not automatically clipped by an under-confident pair head at
inference. The weights ramp up after an initial supervised training period
(Section~\ref{sec:method:loss}), reducing reliance on early, poorly fitted
reference scores without guaranteeing that later references are calibrated.

The cost of this flexibility is that neither exact score ordering nor
thresholded prediction consistency is guaranteed. A low violation rate is not
by itself evidence of better prediction or calibration, and the detached
reference branch is not frozen. Shared encoders, experts and transfer paths
can still carry conflicting gradients. The benefit to be tested is fewer
label-applicable score gaps with acceptable ranking, calibration and classwise
error rates, rather than satisfaction of a universal probability hierarchy.
Numerical comparisons with no consistency, symmetric and label-agnostic
penalties, and hard projection remain \NA{} until those experiments are
reported.

\paragraph{Complementarity with MoE, transfer and auxiliary supervision.}
MoE decides \emph{which interaction transformations to share} across relations.
The PT-to-PHT gate decides \emph{which recognition features to use in the
triple context} and exposes their extractor to triple supervision. Consistency
instead decides \emph{which score to correct when a label-applicable gap
occurs}. Auxiliary pair BCE trains each pair score against its binary label on
triple rows; consistency additionally compares that score with another head's
prediction. The three mechanisms can affect one another through shared
parameters, but none makes the others redundant. In particular, improving
both tasks through feature transfer does not establish output consistency,
and reducing score gaps does not establish that expert routing or feature
transfer is effective. Table~\ref{tab:mechanism-roles} separates these roles.

\subsection{Training objective}\label{sec:method:loss}

The objective sums per-task supervision, auxiliary supervision available only
on triple rows, the two consistency terms under a delayed ramp, and an optional
routing regulariser. Let $\mathrm{BCE}_r$ be binary cross-entropy over the rows
of the current batch with a valid label for relation $r$, computed as a plain
mean with no class re-weighting, and let masked labels never act as negatives.
For a batch of native task $r$ at epoch $e$, let $I_r$ be one for a
triple batch and zero otherwise. The per-batch objective is
\begin{equation}
\begin{aligned}
\mathcal{L}_r \;=\;{} & w_r\,\mathrm{BCE}_r
+\lambda_{\mathrm{bal}}\mathcal{L}_{\mathrm{bal}}\\
&+ I_r\Bigl[\lambda^{\mathrm{aux}}_{\PHm}\mathrm{BCE}_{\PHm}
+\lambda^{\mathrm{aux}}_{\PTm}\mathrm{BCE}_{\PTm}\\
&\qquad +\gamma(e)\bigl(\lambda_{\mathrm{H}}\mathcal{L}_{\mathrm{H}}
+\lambda_{\mathrm{T}}\mathcal{L}_{\mathrm{T}}\bigr)\Bigr],
\end{aligned}
\label{eq:total}
\end{equation}
where only the native task receives its $w_r$-weighted BCE in that batch;
pair BCE on triple rows enters through the auxiliary terms, without being
counted again as native pair supervision. In joint mode, the three per-batch
losses accumulate before an optimiser step.
The auxiliary terms use the \PH{} and \PT{} columns that triple training rows
carry, so they are projected binary supervision on triple rows rather than
additional native data; Section~\ref{sec:method:data} describes the audit that
keeps this projection from leaking into the binary held-out sets. The ramp
\begin{equation}
\gamma(e)=\mathrm{clip}\!\left(\frac{e-E_{\mathrm{warm}}}
{E_{\mathrm{ramp}}},\,0,\,1\right)
\label{eq:ramp}
\end{equation}
holds the consistency terms at zero for the first $E_{\mathrm{warm}}$ epochs
and reaches full weight after a further $E_{\mathrm{ramp}}$, so the heads
are trained against their labels before score coupling receives its full
weight. This schedule does not guarantee coherence or calibrated references. The
routing regulariser $\mathcal{L}_{\mathrm{bal}}=\sum_j(\bar g_{r,j}-\tfrac14)^2$
pulls the batch-mean gate of the current task's router toward uniform. We set
$\lambda_{\mathrm{bal}}=0$ in all reported configurations, so routing
uniformity is \emph{not} an active mechanism behind our results, and we report
the gate statistics in Section~\ref{sec:results:routing} as a diagnostic only.
Table~\ref{tab:hparams} in Appendix~\ref{sec:app:impl} lists all weights.

\begin{algorithm}[!t]
\caption{One \method{} optimisation step in joint mode.}\label{alg:train}
\begin{algorithmic}[1]
\Require batches $B_{\PHm},B_{\PTm},B_{\PHTm}$; epoch $e$; temperature
$\tau(e)$; ramp $\gamma(e)$
\Ensure updated downstream parameters $\theta$
\State \texttt{zero\_grad}($\theta$)
\For{$r \in (\PHm,\PTm,\PHTm)$}
  \State $O \gets \mathrm{forward}(B_r, \tau(e))$
         \Comment{all three logits when $r=\PHTm$}
  \State $\mathcal{L}_r \gets$ Eq.~\eqref{eq:total} on $O$, $\gamma(e)$
  \State \textbf{abort} if $\mathcal{L}_r$ is not finite
  \State $\mathcal{L}_r.\mathrm{backward}()$ \Comment{gradients accumulate}
\EndFor
\State clip $\|\nabla_\theta\|_2$ to $1.0$; one AdamW step; advance schedule
\end{algorithmic}
\end{algorithm}

Two training modes realise Eq.~\eqref{eq:total} at different granularities, and
we report them separately because they weight the tasks differently. In
\textbf{joint} mode (Algorithm~\ref{alg:train}) one microbatch of each task is
forwarded and back-propagated before a single optimiser step, so the three
tasks receive equal numbers of updates and the shorter datasets are recycled
within an epoch. In \textbf{sequential} mode each round makes one full pass over
\PH{}, then \PT{}, then \PHT{}, with an optimiser step per batch, so update
counts follow the natural dataset sizes and the triple task is always last in a
round. The second arrangement is the more realistic use of the available data
but carries a recency bias toward the task seen last, which we flag when
comparing the two in Section~\ref{sec:results:modes}. Validation and early
stopping occur only after all three tasks have been visited.

A single forward pass yields all three logits on triple batches, which is what
makes Eq.~\eqref{eq:total} computable without extra passes. The same pass
therefore supplies the triple prediction, the two auxiliary predictions and the
four consistency operands, so the auxiliary and consistency terms cost no
additional encoder or interaction computation.

\subsection{Data and evaluation protocol}\label{sec:method:data}

\paragraph{Splits.}
We split at the level of entities rather than rows, so that no peptide, allele
or receptor seen in training reappears at evaluation. Entities are allocated
70\%/15\%/15\% under a fixed seed, task rows are derived from that allocation
afterwards, and rows whose unseen entities span more than one type are held out
of the single-entity benchmarks to keep each scenario interpretable. Peptides
are restricted to 8--15 residues. Table~\ref{tab:data} gives the resulting
sizes and Table~\ref{tab:scenarios} the evaluation scenarios.

\begin{table}[!t]
\caption{Data composition after length filtering and class balancing.%
\label{tab:data}}
\begin{tabular*}{\columnwidth}{@{\extracolsep\fill}lrrr@{\extracolsep\fill}}
\toprule
Task & Train & Validation & Test\\
\midrule
\PH{}  & 291\,052 & 168\,680 & 110\,576\\
\PT{}  & 664\,952 & 307\,670 & 298\,034\\
\PHT{} &  72\,192 &  32\,056 &  42\,146\\
\botrule
\end{tabular*}
\begin{tablenotes}%
\item Entity-level split with a 70\%/15\%/15\% allocation under a fixed seed.
\item \PH{} and \PT{} are balanced to a 1:1 positive-to-negative ratio in all
three partitions.
\item The \PHT{} training set is balanced across the four label classes at
18\,048 rows each, so its aggregated positive-to-negative ratio is 1:3;
\PHT{} validation and test are balanced to 1:1 in aggregate.
\item 5\,694 rows with contradictory labels across tasks were quarantined
rather than relabelled.
\end{tablenotes}
\end{table}

\begin{table}[!t]
\caption{Held-out evaluation scenarios. Each withholds one entity type, so no
entity seen in training reappears at test time.\label{tab:scenarios}}
\begin{tabular*}{\columnwidth}{@{\extracolsep\fill}p{0.30\columnwidth}p{0.32\columnwidth}p{0.30\columnwidth}@{\extracolsep\fill}}
\toprule
Scenario & Withheld entity & Tasks\\
\midrule
Unseen peptide & peptide            & \PH{}, \PT{}, \PHT{}\\
Unseen HLA     & HLA allele         & \PH{}, \PHT{}\\
Unseen TCR     & CDR3$\beta$        & \PT{}, \PHT{}\\
Unseen pair    & peptide--receptor  & \PH{}, \PT{}\\
Unseen triple  & peptide--HLA--TCR  & \PHT{}\\
\botrule
\end{tabular*}
\begin{tablenotes}%
\item Ten benchmark files in total. The unseen-pair and unseen-triple
scenarios are combination splits: their individual entities may have been seen
separately, so they are not strictly cold.
\end{tablenotes}
\end{table}

\paragraph{Audit.}
Because our auxiliary terms draw \PH{} and \PT{} supervision from triple rows,
a dedicated audit runs before every training and evaluation action and aborts
on any violation. It validates column names and split tags, filters sequences
that break the length bounds or the canonical 20-residue alphabet while logging
every removal with its source row, rejects any triple row whose class code is
outside the four observed codes or disagrees with its individual labels, drops
exact duplicates, and raises on two rows that share an entity key but carry
different labels. It then checks every pair of splits for overlapping entity
keys. Its decisive check reconstructs the set of binary pairs the model will be
supervised on --- native \PH{}/\PT{} training rows \emph{plus} the \PH{}/\PT{}
columns of the triple training file --- and scans the binary validation, test
and benchmark files for any of them, failing on the first hit. Nothing is
relabelled, dropped or re-split automatically, and a failing corpus simply
cannot be trained on. On the uncleaned corpus this check does fire: five
peptide--receptor pairs projected from triple training rows reappear across ten
rows of the \PT{} test set and the \PT{} unseen-receptor benchmark, with no
\PH{} leakage at all. All reported models are therefore trained on an
explicitly cleaned copy whose fingerprint is recorded in the released
configuration.

\paragraph{Metrics.}
We report threshold-free and threshold-dependent metrics separately, because
the latter depend on a decision rule chosen on validation data. AUROC and AUPRC
are threshold-free. Accuracy, MCC, F1, precision, recall, sensitivity and
specificity use a per-task threshold selected on validation by maximising MCC
and then frozen for all test and benchmark evaluation. We additionally report
negative log-likelihood, Brier score and expected calibration error, since the
consistency objective acts on probabilities and could in principle improve
ordering while degrading calibration. Model selection during training uses the
mean of the three validation AUPRC values.

\subsection{Implementation details}\label{sec:method:impl}

All downstream modules are randomly initialised and trained from scratch; no
checkpoint from an earlier model is loaded, and the checkpoint loader refuses
weights from any other architecture. We use $\dmodel=256$, three encoder layers
with 8 heads and feed-forward width 1024, expert inner width 512, router and
gate inner width 128, and dropout 0.1. Optimisation uses AdamW with weight
decay 0.01 and two learning-rate groups: $2\times10^{-4}$ for all modules and
$1\times10^{-4}$ for the three expert routers and the transfer-gate network,
which we slow down because they control mixing rather than representation.
The schedule warms up linearly over the first 5\% of steps and then decays
linearly; gradients are clipped to global norm 1.0; computation uses
\texttt{bfloat16} autocast on GPU. The routing temperature anneals linearly
from 2.0 to 1.0 over the first five epochs and is then held at 1.0. The
consistency ramp of Eq.~\eqref{eq:ramp} uses $E_{\mathrm{warm}}=3$ and
$E_{\mathrm{ramp}}=5$, so the consistency terms are inactive for three epochs
and reach full weight at epoch eight, while the auxiliary terms are active from
the first epoch. Every configuration, including every ablation variant, is trained three times
under seeds 42, 3407 and 2026, and all reported numbers are means over those
three runs with their standard deviations; the two exceptions are the
post-hoc rows marked in Tables~\ref{tab:ablation}
and~\ref{tab:ablation2}, which involve no training, and the
sensitivity rows of Table~\ref{tab:app:sensitivity}, which are reported at one
seed and labelled as such. Training runs for at most 100 epochs with early stopping after
10 epochs without improvement, where an improvement must exceed
$1\times10^{-4}$. Each task contributes microbatches of 32 rows and evaluation
uses batches of 256. Appendix~\ref{sec:app:impl} lists every hyper-parameter.
\TODO{State the seeds actually run. The implementation takes one seed per
invocation and has no multi-seed runner, so if the tables report a spread over
seeds, say how many independent runs produced it; otherwise report single-seed
results and label them as such.}

%%%%%%%%%%%%%%%%%%%%%%%%%%%%%%%%%%%%%%%%%%%%%%%%%%%%%%%%%%%%%%%%%%%%%%%%%%%%%%
\section{Results}\label{sec:results}
%%%%%%%%%%%%%%%%%%%%%%%%%%%%%%%%%%%%%%%%%%%%%%%%%%%%%%%%%%%%%%%%%%%%%%%%%%%%%%

% Evidence status: the authors confirm that the transfer ablation improves
% both PT and PHT. Numerical results remain on the server and are not yet
% incorporated here. Do not infer that the other listed controls are complete.
\TODO{The qualitative transfer-ablation conclusion is included below.
Numerical values remain \NA{} until the server-side results are incorporated.
Add the actual metrics, evaluation splits, seed counts and training settings
from the run records; distinguish completed comparisons from proposed
mechanistic controls.}

\subsection{Experimental setup}\label{sec:results:setup}

\TODO{Baselines. Nothing in the current code wraps an external predictor, so
state explicitly which baselines were run and how. For the single-relation
families of Sections~\ref{sec:related:ph}--\ref{sec:related:pt}, note that they
score only one relation and leave the other columns blank. For the triple
column, include the two heuristic combinations computable from our own
per-row outputs, $s_{\min}=\min(\hat y_{\PHm},\hat y_{\PTm})$ and
$s_{\mathrm{prod}}=\hat y_{\PHm}\hat y_{\PTm}$, which are the fairest
available controls for whether the triple branch is needed at all.}

\subsection{Comparison with existing predictors}\label{sec:results:main}

\begin{table*}[!t]
\caption{Comparison on the held-out test sets of Table~\ref{tab:data}. Arrows
give the metric direction. Single-relation baselines are evaluated only on the
relations they define.\label{tab:main}}
\tabcolsep=0pt%
\begin{tabular*}{\textwidth}{@{\extracolsep{\fill}}lcccccc@{\extracolsep{\fill}}}
\toprule
& \multicolumn{2}{@{}c@{}}{\PH{}} & \multicolumn{2}{@{}c@{}}{\PT{}}
& \multicolumn{2}{@{}c@{}}{\PHT{}}\\
\cline{2-3}\cline{4-5}\cline{6-7}
Method & AUROC $\uparrow$ & AUPRC $\uparrow$ & AUROC $\uparrow$
& AUPRC $\uparrow$ & AUROC $\uparrow$ & AUPRC $\uparrow$\\
\midrule
\multicolumn{7}{@{}l}{\textit{Single-relation predictors}}\\
NetMHCpan-4.1  & \NA & \NA &     &     &     &    \\
MHCflurry 2.0  & \NA & \NA &     &     &     &    \\
TransPHLA      & \NA & \NA &     &     &     &    \\
ERGO           &     &     & \NA & \NA &     &    \\
pMTnet         &     &     &     &     & \NA & \NA\\
TITAN          &     &     & \NA & \NA &     &    \\
\midrule
\multicolumn{7}{@{}l}{\textit{\PH{}$\rightarrow$\PHT{} transfer models}}\\
% Use the math symbol: OUP redefines \dag as a text-mode superscript.
Fusion-pMT$^{\dagger}$     &     &     &     &     & \NA & \NA\\
TranspMHC / TransTCR    & \NA & \NA &     &     & \NA & \NA\\
StriMap$^{\ddagger}$       & \NA & \NA &     &     & \NA & \NA\\
\midrule
\multicolumn{7}{@{}l}{\textit{Heuristic triple combinations of our own pair heads}}\\
$s_{\min}$             &     &     &     &     & \NA & \NA\\
$s_{\mathrm{prod}}$    &     &     &     &     & \NA & \NA\\
\midrule
\multicolumn{7}{@{}l}{\textit{Unified multi-task models}}\\
UnifyImmun              & \NA & \NA & \NA & \NA &     &    \\
UniPMT                  & \NA & \NA & \NA & \NA & \NA & \NA\\
\method{} (joint)       & \NA & \NA & \NA & \NA & \NA & \NA\\
\method{} (sequential)  & \NA & \NA & \NA & \NA & \NA & \NA\\
\botrule
\end{tabular*}
\begin{tablenotes}%
\item \NA{} marks a result that is pending. A blank cell marks a relation the
method does not explicitly train a head for in the described setting. None of
the transfer models has a peptide--TCR branch; UnifyImmun's specified loss
supervises two pairwise heads, so a derived triple score would require a
separately documented evaluation rule. Fusion-pMT's finetuned model reports
only the triple task.
\item $^{\dagger}$No public implementation at the time of writing (the authors
state that code will be released upon publication), so this row may not be
obtainable. $^{\ddagger}$Requires paired receptor $\alpha$/$\beta$ chains and
predicted structures, which our corpus does not provide; evaluating it on our
splits therefore needs a structure-prediction step and a single-chain
restriction that departs from its published setting.
\item Both metrics are threshold-free, so no decision threshold enters this
table.
\item \TODO{State for each baseline whether it was retrained on our splits or
run as a released model. A released model was fitted on its own corpus, which
may overlap our test entities, so a number obtained that way measures the
published tool rather than the architecture.}
\end{tablenotes}
\end{table*}

\TODO{Prose for Table~\ref{tab:main}, once filled. State the comparison
protocol caveat plainly: external predictors were trained on their own corpora,
so a difference in these columns mixes modelling and training-data effects.}

\subsection{Generalisation to unseen entities}\label{sec:results:unseen}

\begin{table*}[!t]
\caption{Performance on the ten held-out benchmarks of
Table~\ref{tab:scenarios}. Thresholds are the validation MCC-optimal values,
frozen before evaluation.\label{tab:unseen}}
\tabcolsep=0pt%
\begin{tabular*}{\textwidth}{@{\extracolsep{\fill}}llccccc@{\extracolsep{\fill}}}
\toprule
Task & Benchmark & $N$ & AUROC $\uparrow$ & AUPRC $\uparrow$ & MCC $\uparrow$
& F1 $\uparrow$\\
\midrule
\PH{}  & unseen peptide & \NA & \NA & \NA & \NA & \NA\\
\PH{}  & unseen HLA     & \NA & \NA & \NA & \NA & \NA\\
\PH{}  & unseen pair    & \NA & \NA & \NA & \NA & \NA\\
\midrule
\PT{}  & unseen peptide & \NA & \NA & \NA & \NA & \NA\\
\PT{}  & unseen TCR     & \NA & \NA & \NA & \NA & \NA\\
\PT{}  & unseen pair    & \NA & \NA & \NA & \NA & \NA\\
\midrule
\PHT{} & unseen peptide & \NA & \NA & \NA & \NA & \NA\\
\PHT{} & unseen HLA     & \NA & \NA & \NA & \NA & \NA\\
\PHT{} & unseen TCR     & \NA & \NA & \NA & \NA & \NA\\
\PHT{} & unseen triple  & \NA & \NA & \NA & \NA & \NA\\
\botrule
\end{tabular*}
\begin{tablenotes}%
\item Unseen-pair and unseen-triple are combination splits and are not
strictly cold at the level of individual entities.
\item A small number of receptors in the two unseen-TCR benchmarks appear in
some task's training pool --- five in the \PT{} file and four in the \PHT{}
file on the raw corpus --- so strict cross-task cold-receptor performance
cannot be claimed from those two files. \TODO{replace these two counts with the
\texttt{cold\_audit} figures emitted by the reported run, since they depend on
the data copy.}
\end{tablenotes}
\end{table*}

\begin{figure}[!t]
\centering
\figph{5cm}{Figure placeholder. ROC and precision--recall curves per task,
one panel per relation, with one line per held-out scenario.}
\caption{\TODO{ROC/PR caption}}\label{fig:roc}
\figalttext[ROC and precision-recall curves per task]{\TODO{alt text}}%
\end{figure}

\TODO{Prose for Table~\ref{tab:unseen} and Figure~\ref{fig:roc}.}

\subsection{Behaviour across the four label classes}\label{sec:results:classes}

\begin{table}[!t]
\caption{Per-class behaviour on the \PHT{} benchmarks at the frozen threshold:
recall on the positive class and false-positive rate on each negative class.%
\label{tab:classes}}
\begin{tabular*}{\columnwidth}{@{\extracolsep\fill}lcccc@{\extracolsep\fill}}
\toprule
& \texttt{111} & \texttt{000} & \texttt{100} & \texttt{010}\\
Benchmark & recall $\uparrow$ & FPR $\downarrow$ & FPR $\downarrow$
& FPR $\downarrow$\\
\midrule
unseen peptide & \NA & \NA & \NA & \NA\\
unseen HLA     & \NA & \NA & \NA & \NA\\
unseen TCR     & \NA & \NA & \NA & \NA\\
unseen triple  & \NA & \NA & \NA & \NA\\
\botrule
\end{tabular*}
\begin{tablenotes}%
\item \texttt{100} and \texttt{010} test false positives when one pair label
is negative. They can expose reliance on only one pair, but cannot distinguish
triple reasoning from conjunction of both pair labels. Report class counts
with each rate; small denominators limit interpretation.
\end{tablenotes}
\end{table}

\TODO{Compare class-specific false positives with $s_{\min}$ and
$s_{\mathrm{prod}}$, alongside recall on \texttt{111}. Neither these classes
nor a gain over estimated pair scores substitutes for verified \texttt{110}
examples.}

\subsection{Effect of the consistency objective}\label{sec:results:consistency}

\begin{table}[!t]
\caption{Score-gap diagnostics for the full consistency objective on the \PHT{}
test set. A violation is $\hat y_{\PHTm}-\hat y_{\text{pair}}-m>0$, a
numerical diagnostic rather than a biological contradiction.
\label{tab:consistency-results}}
\begin{tabular*}{\columnwidth}{@{\extracolsep\fill}llcc@{\extracolsep\fill}}
\toprule
Coupling & Scope & Violation rate $\downarrow$ & Mean excess $\downarrow$\\
\midrule
$\mathcal{L}_{\mathrm{H}}$ & all        & \NA & \NA\\
$\mathcal{L}_{\mathrm{H}}$ & applicable & \NA & \NA\\
$\mathcal{L}_{\mathrm{T}}$ & all        & \NA & \NA\\
$\mathcal{L}_{\mathrm{T}}$ & applicable & \NA & \NA\\
\botrule
\end{tabular*}
\begin{tablenotes}%
\item ``Applicable'' uses complete-label classes \texttt{111}/\texttt{000}/\texttt{010}
for the PH coupling and \texttt{111}/\texttt{000}/\texttt{100} for the PT
coupling, including rows whose gaps are within the margin. ``All'' uses every
evaluated triple. Report denominators; mean excess is the mean of the positive
part of the gap, including zeros, over the stated scope. Repeat this report
for the comparator rows in Table~\ref{tab:ablation2}.
\item Report AUROC, AUPRC, Brier and ECE alongside these rows, so that a drop
in violations is not mistaken for a globally depressed triple score.
\end{tablenotes}
\end{table}

The consistency objective is evaluated as an additional training preference,
not as enforcement of a probability law. Its results are currently \NA{};
the confirmed PT/PHT transfer gains do not establish a consistency benefit.
For each method, report both the frequency and magnitude of score gaps on the
eligible label groups. The all-row statistic is a separate diagnostic: some
rows are deliberately outside a coupling's training mask, and differences
between conditioning inputs can justify score-order differences.

Compare the full objective with no consistency, a symmetric masked hinge,
label-agnostic coupling, and the hard projection of
Eq.~\eqref{eq:hard-projection}. Keep the MoE, transfer path and auxiliary BCE
unchanged so that this comparison addresses score coupling. Report AUROC and
AUPRC for all three tasks, Brier score and ECE, recall on \texttt{111}, and
triple false-positive rates on the negative classes. Lower violations alone
can result from suppressing all triple scores. Hard projection achieves zero
positive excess by construction, so predictive and calibration metrics are
necessary to assess its cost or benefit; it cannot be judged by violation rate
alone. With task-specific decision thresholds, numerical ordering also need
not ensure agreement between thresholded decisions.

\TODO{Populate the comparison with \NA{} replaced by verified measurements.
Distinguish measured score-gap reduction from changes in ranking, calibration
and classwise errors; do not claim an advantage over hard projection before
those comparisons are available.}

\subsection{Ablation study}\label{sec:results:ablation}

We organise the ablation around a single rule: every qualifier in the
contribution list of Section~\ref{sec:intro} gets one variant that removes
exactly that qualifier and nothing else. Table~\ref{tab:ablation} covers the
sharing decomposition and its placement, together with the controls that ask
whether joint supervision and an explicit triple head improve prediction
on the observed support (contribution~1);
Table~\ref{tab:ablation2} covers the transfer and the consistency objective.
Table~\ref{tab:app:sensitivity} in Appendix~\ref{sec:app:results} reports the
design choices that are not claim-bearing.

\begin{table*}[!t]
\caption{Ablation of the shared--private decomposition and its placement
(contribution~2), with controls for the predictive value of the unified
three-relation setting (contribution~1).\label{tab:ablation}}
\tabcolsep=0pt%
\begin{tabular*}{\textwidth}{@{\extracolsep{\fill}}lccccc@{\extracolsep{\fill}}}
\toprule
& \PH{} & \PT{} & \PHT{} & \PHT{} on \texttt{100}
& \PHT{} on \texttt{010}\\
Variant & AUPRC $\uparrow$ & AUPRC $\uparrow$ & AUPRC $\uparrow$
& FPR $\downarrow$ & FPR $\downarrow$\\
\midrule
Full model                             & \NA & \NA & \NA & \NA & \NA\\
\midrule
\multicolumn{6}{@{}l}{\textit{Contribution 1: what does joint three-relation learning add?}}\\
\PH{} only, single relation            & \NA &     &     &     &     \\
\PT{} only, single relation            &     & \NA &     &     &     \\
\PHT{} only, no pair heads             &     &     & \NA & \NA & \NA\\
$s_{\min}$ from our own pair heads$^{\ddagger}$
                                       &     &     & \NA & \NA & \NA\\
$s_{\mathrm{prod}}$ from our own pair heads$^{\ddagger}$
                                       &     &     & \NA & \NA & \NA\\
\midrule
\multicolumn{6}{@{}l}{\textit{Contribution 2: is the decomposition right?}}\\
Plain FFN sublayer, matched parameters & \NA & \NA & \NA & \NA & \NA\\
All-shared experts ($4+0$)             & \NA & \NA & \NA & \NA & \NA\\
Private-only experts ($0+4$)           & \NA & \NA & \NA & \NA & \NA\\
All-shared, matched total parameters   & \NA & \NA & \NA & \NA & \NA\\
Private-only, matched total parameters & \NA & \NA & \NA & \NA & \NA\\
Uniform (non-learned) gate             & \NA & \NA & \NA & \NA & \NA\\
Gate on entity summaries only          & \NA & \NA & \NA & \NA & \NA\\
Gate on peptide summary only           & \NA & \NA & \NA & \NA & \NA\\
\midrule
\multicolumn{6}{@{}l}{\textit{Contribution 2: is the placement right?}}\\
Experts after pooling                  & \NA & \NA & \NA & \NA & \NA\\
Triple reads pre-expert $U_{\PHm}$     & \NA & \NA & \NA & \NA & \NA\\
Shared pool over $\{\PHm,\PTm\}$ only  & \NA & \NA & \NA & \NA & \NA\\
Shared pool over $\{\PHm,\PHTm\}$ only & \NA & \NA & \NA & \NA & \NA\\
Shared pool over $\{\PTm,\PHTm\}$ only & \NA & \NA & \NA & \NA & \NA\\
\botrule
\end{tabular*}
\begin{tablenotes}%
\item Every cell is the mean over seeds 42, 3407 and 2026 with its standard
deviation; a single-seed row is never comparable to another at the effect sizes
at stake here. A blank cell marks a relation the variant does not define.
\item The last two columns report triple false-positive rates on
\texttt{100} and \texttt{010}, using each model's threshold selected on the
full validation set and frozen before testing. Both classes have negative
\PHT{} labels, so AUPRC on either one against \texttt{000} would not be a
meaningful classification measure. Report denominators with the rates: the
raw test split contains only four \texttt{100} rows and 331 \texttt{010} rows.
These columns test rejection when one pair condition fails, not recognition
beyond conjunction. Interpret them alongside positive-class recall to avoid
rewarding uniformly low triple scores.
\item $^{\ddagger}$Computed post hoc from the full model's own \PH{} and \PT{}
outputs, so these two rows involve no training. Together with the
\PHT{}-only row they are the fair answer to whether the triple branch adds
anything beyond combining its own pair predictions, which
Section~\ref{sec:discussion} identifies as the sharpest question this corpus
permits.
\item The three single-relation rows cannot use the model-selection rule of
Section~\ref{sec:method:data}, which averages three validation AUPRC
values; they select on their own relation's validation AUPRC. All other rows
use the unchanged rule, and no row changes the epoch budget or the
early-stopping criterion.
\item The $4+0$ and $0+4$ controls keep four active experts per branch,
not the total parameter count: the full model has eight unique experts, the
all-shared model four and the private-only model twelve. The additional
parameter-matched rows adjust expert widths and must report actual trainable
parameter counts. Only those matched controls support attribution beyond
added capacity.
\item ``Gate on entity summaries only'' drops the interaction term
$\bar U_r$ from Eq.~\eqref{eq:context} and leaves the entity means, which is
the gate-input condition of the multi-task originals
(Section~\ref{sec:related:moe}); ``gate on peptide summary only'' additionally
drops the receptor, and so also removes the gate's access to the relation.
\item The three restricted-pool rows test sharing across relations with
different input compositions. If gains occur only for $\{\PHm,\PHTm\}$,
which both receive HLA, evidence for sharing between the HLA-conditioned and
HLA-unspecified relations is limited. Such a result does not by itself reveal
which biological features the experts represent.
\end{tablenotes}
\end{table*}

\TODO{Prose for Table~\ref{tab:ablation}. Evaluate contribution~2 using
both task performance and interference diagnostics. On fixed complete-label
batches, compute pairwise cosine similarity between gradients of the native
supervised task losses on the same shared parameter subset; report encoders
and shared experts separately. A negative cosine indicates a locally
conflicting update direction, not inevitable performance degradation. Compare
with matched sharing controls and report joint-minus-single-task performance.
Until these measurements exist, describe mitigation as the design aim rather
than an achieved result.}

\begin{table*}[!t]
\caption{Ablation of the gated residue-level transfer and supervision feedback
(contribution~3), the
label-aware consistency objective (contribution~4) and the auxiliary
supervision. Metric columns and the intended reporting protocol follow
Table~\ref{tab:ablation}.\label{tab:ablation2}}
\tabcolsep=0pt%
\begin{tabular*}{\textwidth}{@{\extracolsep{\fill}}lccccc@{\extracolsep{\fill}}}
\toprule
& \PH{} & \PT{} & \PHT{} & \PHT{} on \texttt{100}
& \PHT{} on \texttt{010}\\
Variant & AUPRC $\uparrow$ & AUPRC $\uparrow$ & AUPRC $\uparrow$
& FPR $\downarrow$ & FPR $\downarrow$\\
\midrule
Full model                             & \NA & \NA & \NA & \NA & \NA\\
\midrule
\multicolumn{6}{@{}l}{\textit{Contribution 3: which property of the transfer
matters?}}\\
w/o gated transfer, path removed$^{\S}$ & \NA & \NA & \NA & \NA & \NA\\
Gated transfer, detached PT features    & \NA & \NA & \NA & \NA & \NA\\
Ungated transfer, $\beta\equiv1$       & \NA & \NA & \NA & \NA & \NA\\
Scalar $\beta$, one per sample         & \NA & \NA & \NA & \NA & \NA\\
$\beta$ without explicit HLA summary   & \NA & \NA & \NA & \NA & \NA\\
Transfer of pre-expert $U_{\PTm}$      & \NA & \NA & \NA & \NA & \NA\\
Pooled-vector fusion instead           & \NA & \NA & \NA & \NA & \NA\\
\midrule
\multicolumn{6}{@{}l}{\textit{Contribution 4: which property of the objective
matters?}}\\
w/o consistency terms                  & \NA & \NA & \NA & \NA & \NA\\
Hierarchy term only, $\lambda_{\mathrm{T}}=0$
                                       & \NA & \NA & \NA & \NA & \NA\\
\PT{}-implication only, $\lambda_{\mathrm{H}}=0$
                                       & \NA & \NA & \NA & \NA & \NA\\
Label-agnostic consistency             & \NA & \NA & \NA & \NA & \NA\\
Symmetric hinge, no stop-gradient      & \NA & \NA & \NA & \NA & \NA\\
Hard $\min$ constraint during training & \NA & \NA & \NA & \NA & \NA\\
Hard $\min$ at inference only$^{\ddagger}$ & \NA & \NA & \NA & \NA & \NA\\
\midrule
\multicolumn{6}{@{}l}{\textit{Auxiliary supervision on triple rows}}\\
w/o auxiliary pair supervision         & \NA & \NA & \NA & \NA & \NA\\
Auxiliary \PH{} only                   & \NA & \NA & \NA & \NA & \NA\\
Auxiliary \PT{} only                   & \NA & \NA & \NA & \NA & \NA\\
\botrule
\end{tabular*}
\begin{tablenotes}%
\item Numerical values are \NA{} pending integration of the server-side
results. The confirmed conclusion concerns the improvement in both \PT{}
and \PHT{} from adding gated transfer. Other rows define diagnostic controls;
their presence does not assert that those experiments have been completed.
Record the actual evaluation protocol and seed counts with the results.
\item Detached transfer uses $\sg[Z_{\PTm}]$ in both the gate input and the
residual projection. It preserves the forward computation at fixed parameters
but blocks both terms of Eq.~\eqref{eq:transfer-feedback}. The \PT{} head
still uses the original, non-detached representation for its own supervision.
Shared encoders and experts may still receive triple gradients through other
paths, so this control isolates feedback through the added connection only.
\item The first block separates the four qualifiers of contribution~3, which a
single row removing the whole path cannot do. $\beta\equiv1$ tests whether the
transfer has to be \emph{gated}; a scalar $\beta$ keeps the gate but removes
its residue resolution, which is what ``residue-level'' means; dropping
$\bar H_H$ from Eq.~\eqref{eq:beta} tests the added value of an explicit HLA
summary. This control still receives HLA information through the triple state,
so it is not HLA-independent. The
pooled-vector row replaces the residue injection with a pooled \PT{} feature
concatenated at the triple classifier. This is a control for retaining the
residue axis, not a reproduction of Fusion-pMT, which already transfers
unpooled \PH{} features.
\item $^{\S}$Disabling the transfer leaves its parameters registered but
unused, so this row is not matched in active parameters. The $\beta\equiv1$
control retains the same feature projection $R$ but removes learned gating;
the gate network is then inactive, which must be reflected in active-parameter
counts. The detached-transfer control retains all forward modules.
\item ``Label-agnostic consistency'' uses the fixed-direction penalty
$\phi(\hat y_{\PHTm},\sg[\hat y_{\PHm}])+
\phi(\hat y_{\PHTm},\sg[\hat y_{\PTm}])$ on every complete-label row,
with the respective $\lambda_{\mathrm{H}}$ and $\lambda_{\mathrm{T}}$
weights, margin and ramp unchanged.
It always corrects the triple downward when a gap occurs, removing both the
class masks and the label-dependent choice of correction target. It does not
isolate masking alone. The symmetric-hinge row retains the original class
masks and weights but removes stop-gradient, isolating the effect of letting
both operands receive direct gradients.
\item The hard-$\min$ controls use $q_{\PHTm}$ from
Eq.~\eqref{eq:hard-projection}, which includes the raw triple score and is
not the pair-only baseline $s_{\min}$. The training variant replaces the raw
triple score by $q_{\PHTm}$ in the native triple BCE and omits soft consistency;
pair and auxiliary supervision remain unchanged. At inference it also uses
$q_{\PHTm}$. $^{\ddagger}$The inference-only variant post-processes the full
model without retraining. These answer different questions: constrained
learning versus clipping an already trained predictor.
\item Separating the two consistency terms matters because they are weighted
independently in Eq.~\eqref{eq:total} and act on different label classes:
$\mathcal{L}_{\mathrm{H}}$ is active on \texttt{111}, \texttt{000} and
\texttt{010}; $\mathcal{L}_{\mathrm{T}}$ on \texttt{111}, \texttt{000} and
\texttt{100}.
\item Report AUROC, Brier and ECE for every row of the last two blocks
alongside AUPRC, so that a reduction in violation rate is not mistaken for an
improvement in ranking or calibration.
\end{tablenotes}
\end{table*}

\paragraph{Gated transfer improves both PT and PHT.}
Adding the HLA-conditioned gated connection improves prediction performance
for both \PT{} and \PHT{} relative to removing that connection. The
improvement is therefore not confined to the receiving triple task: the
source pair task also benefits when its representation participates in the
triple computation. The \PHT{} gain supports using independently supervised
recognition features alongside \PH{} features, while the \PT{} gain is
consistent with the additional representation-training path described in
Eq.~\eqref{eq:transfer-feedback}. Numerical results are \NA{} in
Table~\ref{tab:ablation2} pending transfer of the server-side measurements.

The comparison establishes the benefit of the complete gated connection; it
does not isolate each mechanism. Comparing the full connection with detached
transfer tests the contribution of its backward feedback, while comparison
with $\beta\equiv1$ tests conditional gating against unconditional projected
feature transfer. Removing the explicit HLA summary from the gate tests its
incremental value beyond the HLA information already present in the triple
state. These controls should keep task data, auxiliary pair losses,
consistency terms, training budgets and model-selection rules unchanged.
In particular, the effect of the added feature path should not be conflated
with \PT{} auxiliary BCE or label-aware consistency, which provide separate
supervision routes.

\TODO{Insert effect sizes and run metadata for the completed transfer
comparison. Report the diagnostic controls only after their results are
available; do not infer significance or a feedback-specific gain from the
qualitative full-connection ablation. Complete the consistency ablation
discussion separately.}

\TODO{Implementation status. Check which diagnostic variants are implemented
in the server-side experiment code before reporting them. Required controls
include detached PT features at both fusion inputs. In the current local
design, expert counts
and the shared pool's membership are fixed in the model, there is no
gate-override or gate-context switch, no post-pooling expert variant, no
pre-expert triple path, no $\beta$ override, no label-agnostic or hard-$\min$
consistency mode, and no single-relation training mode (the three task weights
are validated as strictly positive). Add them as default-off switches on the
\texttt{stage\_final} subclasses so that no finished stage changes behaviour.
Each switch belongs to the resume contract, so every variant needs its own
configuration file and trains from scratch.}

\subsection{Routing and transfer-gate analysis}\label{sec:results:routing}

\begin{table}[!t]
\caption{Learned allocation between shared and branch-private experts, and
transfer-gate statistics, on the \PHT{} test set. $\alpha_r$ is the shared
fraction $g_{r,1}+g_{r,2}$.\label{tab:routing}}
\begin{tabular*}{\columnwidth}{@{\extracolsep\fill}lccc@{\extracolsep\fill}}
\toprule
Relation & $\alpha_r$ & Gate entropy & \TODO{units}\\
\midrule
\PH{}  & \NA & \NA & \NA\\
\PT{}  & \NA & \NA & \NA\\
\PHT{} & \NA & \NA & \NA\\
\midrule
\multicolumn{4}{@{}l}{\textit{Transfer gate $\beta$ over real peptide residues}}\\
mean / s.d.          & \NA & \NA & \\
fraction $<0.05$     & \NA &     & \\
fraction $>0.95$     & \NA &     & \\
\botrule
\end{tabular*}
\begin{tablenotes}%
\item The gate bias starts at $-2$, corresponding to $\sigma(-2)\approx0.12$
before the input-dependent logit contribution. Compare the measured initial
and final gate distributions to assess how training changes transfer.
\item The routing regulariser is disabled in all reported runs, so $\alpha_r$
reflects learning and not a uniformity penalty.
\end{tablenotes}
\end{table}

\begin{figure}[!t]
\centering
\figph{4.6cm}{Figure placeholder. Left: shared fraction $\alpha_r$ per
relation over training epochs. Right: distribution of the per-residue transfer
gate $\beta$, by peptide position and by HLA supertype.}
\caption{\TODO{routing and gate caption}}\label{fig:gates}
\figalttext[Expert routing and transfer-gate statistics]{\TODO{alt text}}%
\end{figure}

\TODO{Report the empirical tests of contribution~2: variation in $\alpha_r$
between relations and across samples, and performance under matched shared-only,
private-only and fixed-gate controls. These statistics describe allocation of
transformations, not biological contributions. If routing is nearly constant,
state that sample-adaptive allocation is unsupported; distinguish a useful
fixed mixture from a pure capacity gain using the controls, rather than
inferring either from gate values alone. Analyse $\beta$ separately, as a
per-residue transfer gate, and do not interpret it as a physical contact map
without independent structural evidence.}

\subsection{Joint versus sequential training}\label{sec:results:modes}

\TODO{Prose and a small table or figure. State that the two modes give the
tasks different numbers of updates, and that the triple task is last in every
sequential round, so a sequential advantage on \PHT{} may partly reflect
recency.}

\subsection{Computational cost}\label{sec:results:cost}

\TODO{Report trainable parameter count, expert parameter count, peak GPU
memory, optimiser steps and samples per second from \texttt{history.json}.
State that the language models are excluded from these numbers because they
are not in the training graph, and report the one-off embedding precomputation
cost separately.}

%%%%%%%%%%%%%%%%%%%%%%%%%%%%%%%%%%%%%%%%%%%%%%%%%%%%%%%%%%%%%%%%%%%%%%%%%%%%%%
\section{Discussion}\label{sec:discussion}
%%%%%%%%%%%%%%%%%%%%%%%%%%%%%%%%%%%%%%%%%%%%%%%%%%%%%%%%%%%%%%%%%%%%%%%%%%%%%%

The first two contributions address different questions. The unified task
formulation specifies which predictions are available and which molecular
inputs each can use. The MoE specifies how those tasks share interaction
capacity. UniPMT already establishes the value of considering all three
relations in one framework; our design instead gives each relation a direct
sequence branch and makes the \PT{} representation available for
HLA-conditioned residue-level transfer. UnifyImmun and Fusion-pMT motivate
useful sharing and transfer paths, but their published training schemes do not
establish the behaviour of three objectives kept active together. Our
shared--private design addresses that joint setting without assuming that
those earlier architectures have been shown to fail.

The transfer ablation shows that the added gate improves both \PT{} and
\PHT{}, supporting a benefit for the source as well as the receiving task.
For \PHT{}, the connection complements presentation-derived features with
interaction features learned under independent \PT{} supervision, extending
the PH-to-triple transfer used by Fusion-pMT. For \PT{}, the differentiable
path lets the triple objective shape the representation consumed by its own
classifier. The HLA-conditioned gate provides flexibility in using that
representation when the pair and triple labels differ, as in \texttt{010}.
This is a representation-level training interaction: the triple BCE does not
directly train the PT classifier through the transfer edge, and PT inference
does not acquire an HLA input.

These paths offer a mechanistic account of the observed gains, but the
full-connection ablation alone does not separate their contributions. A
detached-transfer control is needed to attribute a gain specifically to
backward feedback, and an ungated control to attribute it to conditional
gating. Moreover, the additional objective need not supply independent
biological information when the labels obey pairwise conjunction; it can
still shape a shared representation usefully. The improvement of both tasks
therefore supports the gated connection without establishing that every
feedback gradient is beneficial or that the missing \texttt{110} distinction
has been learned.

Consistency addresses a third issue beyond shared capacity and feature
transfer: how to correct disagreements among the output scores. A shared
expert can improve features without resolving a pair--triple score gap, and
the transfer gate can improve both tasks without specifying which output
should change when their scores differ. The label-aware objective makes that
choice explicit. For positive triples it uses the triple score to raise a
lagging positive pair; for negative pair--triple labels it uses the pair score
to lower the triple. Unlike a symmetric penalty, each term has only one
active score operand. Unlike hard clipping, it uses labels during training
to choose the correction and leaves inference scores unconstrained.

This flexibility is useful because the three predictions condition on
different molecular inputs. It avoids treating the pair-score bound as a
universal probability law, while still using the annotated relationships as
a training preference. The margin and loss weights allow the model to retain
score differences when demanded by its classification objectives. The price
is that exact hierarchy is not guaranteed, and label-aligned direct
derivatives need not translate into conflict-free shared-parameter updates.
Consistency benefits therefore require their own comparisons: neither the
transfer ablation nor routing statistics establish them. The relevant
question is whether coupling reduces label-applicable score gaps while
preserving or improving prediction and calibration; those numerical results
remain \NA{}.

The rationale for a separate triple task is to avoid equating a regularity in
the available labels with a universal recognition rule. Two positive pair
predictions need not establish recognition under the specified HLA. Learning
a triple decision from all three entities leaves this distinction expressible,
whereas defining that decision as pairwise conjunction rules it out by design.
The triple head thus removes a structural restriction; it does not by itself
correct the preference supported by the observed label distribution.

The absent \texttt{110} class sets the boundary of what the present evaluation
can establish. The \texttt{111} rows fit both an explicit triple predictor and
a conjunction rule; \texttt{100} and \texttt{010} test failure of one pair
condition, which that rule also represents. Improvements over minimum or
product combinations of predicted scores would therefore support predictive
utility on the observed distribution, not prove that the model has learned
HLA-specific exceptions to pairwise positivity. A decisive extension would
collect verified triple negatives with positive pair evidence and a documented
HLA context, rather than equating randomly recombined triples with confirmed
\texttt{110} examples. It would also require updating the current four-class
reader and loss contract before retraining and evaluation.

\TODO{Add the quantitative transfer results and the remaining result-dependent
discussion: comparison with single-task learning, matched expert-sharing
controls, interference diagnostics, and ranking/calibration under consistency.
The confirmed PT/PHT transfer gains do not establish these other outcomes.}

\begin{boxtext}[Limitations]
\begin{unlist}
\item \textbf{Missing evidence for non-conjunctive triple decisions.} No
complete triple-labelled row has class \texttt{110}, and the current reader
and loss reject it. Thus the available labels obey
Eq.~\eqref{eq:observed-conjunction}; the evaluation cannot establish that a
separate triple decision is indispensable. Pair-score composition controls
measure predictive value under this restriction, while verified \texttt{110}
cases and an extended training contract are required to test the broader task.
\item \textbf{Consistency is a soft prior, not a guaranteed hierarchy.}
The class masks and stop-gradient align the penalty's direct correction with
its target label. They neither certify the reference score nor freeze its
branch, and they guarantee neither calibrated probabilities nor consistent
thresholded predictions. Performance relative to hard projection remains
\NA{} pending the dedicated comparison.
\item \textbf{Selective sharing does not isolate gradients.} Shared modules
still receive all task gradients, and the triple loss also reaches pair-private
experts through the transfer paths. The MoE is a capacity-allocation mechanism;
reduced interference and gains over single-task learning require measurement.
\item \textbf{Sparse diagnostic negatives.} The raw test split has only four
\texttt{100} rows, so one error changes its false-positive rate by 25 percentage
points. Report class counts and uncertainty, and avoid drawing conclusions
about triple-specific recognition from this subset.

\item \textbf{The expert mixture is dense and the routing is sample-level.}
All four experts are evaluated for every sample, so we make no efficiency claim
against sparse expert layers, and the gate is computed once per sample, so the
routing is not residue-level even though expert computation is.
\item \textbf{Routing uniformity is not an active mechanism.} The balance
regulariser is available but disabled in all reported runs, so the observed
expert allocation must not be attributed to it.
\item \textbf{Transfer gains and joint-training gains are different comparisons.}
The gated-connection ablation establishes improvement in \PT{} and \PHT{}
within the joint model; it does not establish improvement over training each
task alone. The three
single-relation rows of Table~\ref{tab:ablation} are what licenses any
statement about the unified setting, and they require a training mode that
restricts the active task set, since zeroing a task weight would still let the
other tasks' gradients reach the shared modules. Absent a gain on those rows we
claim only that the three relations can be served by one model, not that they
improve one another.
\item \textbf{Train and evaluation class priors differ.} The triple training
set is balanced across the four label classes, which makes its aggregate
positive-to-negative ratio 1:3, while triple validation and test are balanced
1:1. The loss applies no re-weighting, and this prior shift interacts with the
MCC-selected threshold, so all threshold-dependent metrics should be read with
that in mind.
\item \textbf{Auxiliary supervision is projected, not additional.} The \PH{}
and \PT{} auxiliary terms reuse labels attached to triple rows. The audit of
Section~\ref{sec:method:data} exists because those pairs can otherwise appear
in the binary held-out sets, and on the uncleaned corpus they do; every
reported model is trained on a cleaned copy.
\item \textbf{Mechanisms are not co-timed.} Auxiliary supervision is active
from the first epoch while the consistency terms ramp in from epoch four to
epoch eight, so attributing a gain to one rather than the other requires the
ablations of Table~\ref{tab:ablation} and not the training curves alone.
\item \textbf{No external validation cohort.} All evaluation uses held-out
splits of a single curated corpus. \TODO{state whether an independent
external cohort was obtained; if not, say so here plainly.}
\end{unlist}
\end{boxtext}

%%%%%%%%%%%%%%%%%%%%%%%%%%%%%%%%%%%%%%%%%%%%%%%%%%%%%%%%%%%%%%%%%%%%%%%%%%%%%%
\section{Conclusion}\label{sec:conclusion}
%%%%%%%%%%%%%%%%%%%%%%%%%%%%%%%%%%%%%%%%%%%%%%%%%%%%%%%%%%%%%%%%%%%%%%%%%%%%%%

We presented \method{}, a sequence-based model that jointly serves \PH{},
\PT{} and \PHT{} while respecting their different molecular inputs. The
separate triple head avoids imposing pairwise conjunction as a fixed decision
rule and permits recognition to depend on the specified HLA context. Its
unification differs from UniPMT through direct relation branches, including a
\PT{} representation that does not require averaging over candidate HLA
molecules. Its second contribution adapts shared--private experts to the
interaction blocks, allowing sample-dependent reuse of transformations while
preserving peptide positions for the triple branch. HLA-conditioned \PT{}
transfer complements the PH-derived triple representation with independently
supervised recognition features and returns triple gradients to the PT feature
extractor. The transfer ablation shows improved prediction for both tasks.
Label-aware directional consistency supplies the complementary output-level
objective: it chooses a label-supported correction target for a pair--triple
score gap, rather than enforcing a hard bound on every triple prediction.
MoE allocates transformations, gated transfer exchanges features and their
supervision, and consistency couples scores. Its predictive and calibration
benefits, as well as the individual roles of transfer, feedback and gating,
require their respective controls.
The absence of \texttt{110} supervision remains a concrete limitation:
verified examples of positive pairs with a negative specified triple, together
with an extended data/loss contract, are needed to test recognition beyond
pair-label conjunction.
\TODO{Insert the transfer effect sizes, currently \NA{}, and other
quantitative conclusions supported by the completed experiments.}

%%%%%%%%%%%%%%%%%%%%%%%%%%%%%%%%%%%%%%%%%%%%%%%%%%%%%%%%%%%%%%%%%%%%%%%%%%%%%%
%% Back matter
%%%%%%%%%%%%%%%%%%%%%%%%%%%%%%%%%%%%%%%%%%%%%%%%%%%%%%%%%%%%%%%%%%%%%%%%%%%%%%

\section*{Conflicts of interest}
The authors declare that they have no competing interests.

\section*{Funding}
\TODO{funding statement, with grant numbers}

\section*{Data availability}
\TODO{Data availability statement. Name the repository and identifier for the
curated corpus, and state that the split audit report and the cleaned-copy
fingerprint are released alongside the code so that the splits can be
reproduced exactly.}

\section*{Author contributions statement}
\TODO{contributions by author initials}

\section*{Acknowledgments}
\TODO{acknowledgments}

\bibliographystyle{oup-abbrvnat}
\bibliography{references}

%% Author biographies are required by some OUP journals (for example Briefings
%% in Bioinformatics) and not by Bioinformatics. Delete if not required.
%\begin{biography}{{\color{black!20}\rule{77pt}{57pt}}}{%
%\author{Name.} \TODO{biography}}
%\end{biography}

\begin{appendices}

\section{Hyper-parameters}\label{sec:app:impl}

\begin{table}[!t]
\caption{Complete hyper-parameter settings. Values marked \NA{} are read from
the released configuration file.\label{tab:hparams}}
\begin{tabular*}{\columnwidth}%
{@{\extracolsep\fill}p{0.56\columnwidth}r@{\extracolsep\fill}}
\toprule
Hyper-parameter & Value\\
\midrule
\multicolumn{2}{@{}l}{\textit{Representation}}\\
Peptide/HLA language model (frozen) & ESM-C 300M\\
TCR language model (frozen)         & TCR-BERT\\
Cached embedding dtype              & \texttt{float16}\\
Peptide length range                & 8--15\\
HLA / TCR maximum length            & 34 / 34\\
\midrule
\multicolumn{2}{@{}l}{\textit{Architecture}}\\
Model width $\dmodel$               & 256\\
Encoder layers / heads / FFN width  & 3 / 8 / 1024\\
Shared experts / private experts per relation & 2 / 2\\
Expert inner width                  & 512\\
Router inner width                  & 128\\
Transfer-gate inner width           & 128\\
Transfer-gate bias initialisation   & $-2$\\
Pooling / classifier inner width    & 128 / 128\\
Dropout                             & 0.1\\
\midrule
\multicolumn{2}{@{}l}{\textit{Objective}}\\
Task weights $w_{\PHm},w_{\PTm},w_{\PHTm}$ & 1.0, 1.0, 1.0\\
Auxiliary weights $\lambda^{\mathrm{aux}}_{\PHm},
\lambda^{\mathrm{aux}}_{\PTm}$   & 0.1, 0.1\\
Consistency weights $\lambda_{\mathrm{H}},\lambda_{\mathrm{T}}$ & 0.05, 0.03\\
Consistency margin $m$              & 0.05\\
Ramp $E_{\mathrm{warm}}$ / $E_{\mathrm{ramp}}$ & 3 / 5\\
Balance weight $\lambda_{\mathrm{bal}}$ & 0.0\\
\midrule
\multicolumn{2}{@{}l}{\textit{Optimisation}}\\
Optimiser / weight decay            & AdamW / 0.01\\
Learning rate, modules / gating & $2\times10^{-4}$ / $1\times10^{-4}$\\
Warm-up fraction                    & 0.05\\
Gradient clipping (global norm)     & 1.0\\
Precision                           & \texttt{bfloat16} autocast\\
Routing temperature $\tau$          & $2.0\rightarrow1.0$ over 5 epochs\\
Batch size per task / evaluation    & 32 / 256\\
Maximum epochs / patience           & 100 / 10\\
Minimum improvement $\delta$        & $1\times10^{-4}$\\
Selection metric                    & mean validation AUPRC\\
Seeds                               & 42, 3407, 2026\\
\botrule
\end{tabular*}
\end{table}

\section{Additional results}\label{sec:app:results}

\begin{figure}[!t]
\centering
\figph{4.6cm}{Appendix figure placeholder. Training and validation curves per
task, with the consistency-ramp window and the temperature-annealing window
marked.}
\caption{\TODO{appendix figure caption}}\label{fig:app:curves}
\figalttext[Training and validation curves per task]{\TODO{alt text}}%
\end{figure}

\begin{table}[!t]
\caption{Sensitivity to design choices that no contribution claims. Single
seed (42), since these rows are reported to show that the configuration of
Section~\ref{sec:method:impl} is not a sensitive
one.\label{tab:app:sensitivity}}
\begin{tabular*}{\columnwidth}{@{\extracolsep\fill}p{0.38\columnwidth}ccc@{\extracolsep\fill}}
\toprule
Variant & \PH{} & \PT{} & \PHT{}\\
 & AUPRC $\uparrow$ & AUPRC $\uparrow$ & AUPRC $\uparrow$\\
\midrule
Reported configuration             & \NA & \NA & \NA\\
\midrule
Experts $1+1$                      & \NA & \NA & \NA\\
Experts $4+4$                      & \NA & \NA & \NA\\
No temperature annealing, $\tau\equiv1$ & \NA & \NA & \NA\\
Routing regulariser on, $\lambda_{\mathrm{bal}}>0$ & \NA & \NA & \NA\\
No consistency ramp, $\gamma\equiv1$ & \NA & \NA & \NA\\
Consistency margin $m=0$           & \NA & \NA & \NA\\
Transfer-gate bias init $0$        & \NA & \NA & \NA\\
Single learning rate for routers and gate & \NA & \NA & \NA\\
One encoder layer instead of three & \NA & \NA & \NA\\
\botrule
\end{tabular*}
\begin{tablenotes}%
\item Every row here is reachable by configuration alone, which is why they
are reported together and at one seed. The expert-count rows are the exception
in spirit rather than in cost: they substantiate the statement in
Section~\ref{sec:related:moe} that the two-plus-two split is a setting within a
known scheme rather than a tuned novelty.
\item The $\lambda_{\mathrm{bal}}>0$ row exists so that the routing analysis of
Section~\ref{sec:results:routing} cannot be read as an artefact of a
uniformity penalty that we did not use.
\item The single-learning-rate row tests the stated reason for the second
parameter group in Section~\ref{sec:method:impl}, that routers and the
transfer gate control mixing rather than representation.
\item Joint versus sequential training is a mode rather than an ablation and
is reported in Section~\ref{sec:results:modes}.
\end{tablenotes}
\end{table}

\begin{table}[!t]
\caption{Per-component loss values over training, with their own
denominators.\label{tab:app:loss}}
\begin{tabular*}{\columnwidth}{@{\extracolsep\fill}p{0.43\columnwidth}p{0.20\columnwidth}p{0.25\columnwidth}@{\extracolsep\fill}}
\toprule
Component & Final value & Rows contributing\\
\midrule
\PH{} / \PT{} / \PHT{} supervised & \NA & \NA\\
Auxiliary \PH{} / \PT{}           & \NA & \NA\\
$\mathcal{L}_{\mathrm{H}}$ positive / negative & \NA & \NA\\
$\mathcal{L}_{\mathrm{T}}$ positive / negative & \NA & \NA\\
\botrule
\end{tabular*}
\begin{tablenotes}%
\item Each component is averaged over the rows to which it applies, so the
columns are not comparable as totals.
\end{tablenotes}
\end{table}

\end{appendices}

\end{document}
